% संपूर्ण मराठी ग्रंथातील प्रकरण 26: अनिर्णेयता
% हा संपादनयोग्य स्रोतखंड आहे. संकलनासाठी संपूर्ण स्रोतसंग्रहातील
% अक्षररूपे, पूर्वभाग, चिन्हव्याख्या आणि आकृती-स्रोत आवश्यक आहेत.
% मूळ: Open Logic Project; मजकूर आणि रूपांतर CC BY 4.0.
% यंत्रानुवाद, दुरुस्ती आणि तपासणी: OpenAI Codex —
% GPT-5.6 Sol आणि GPT-6 Sol, दोन्ही Ultra effort.
% दुरुस्ती-पुनरावलोकन: GPT-6.1 Sol, Ultra effort.
\chapter{अनिर्णेयता}\label{tur:und::chap}








\section{प्रस्तावना}\label{tur:und:int:sec}

प्रत्येक फलन, अगदी प्रत्येक अंकगणितीय फलनही, संगणनक्षम
असेलच असे नाही, हे उघड वाटू शकते. अशी फलने खूप आहेत
आणि त्यांचे वर्तन फार गुंतागुंतीचे असू शकते. उदाहरणार्थ,
किरणोत्सर्गी कणांच्या क्षयावरून किंवा इतर अव्यवस्थित अथवा
यादृच्छिक वर्तनावरून परिभाषित केलेली फलने. समजा, एखाद्या
दिलेल्या क्षणापासून आपण एका सेकंदाच्या अंतरांची मोजणी सुरू
करतो आणि त्या आरंभीच्या क्षणानंतरच्या $n$-व्या एक-सेकंदाच्या
अंतरात विश्वातील जितक्या कणांचा क्षय होतो ती संख्या $f(n)$
या फलनाची किंमत म्हणून ठरवतो. हे फलन कधीही संगणित
करता येणार नाही असे वाटण्यासारखे आहे.

मात्र अशी फलने कशी संगणित करावीत याची कल्पना न सुचणे
ही एक गोष्ट आहे; ती संगणनक्षम नाहीत हे प्रत्यक्ष सिद्ध
करणे ही दुसरी गोष्ट आहे. खरे तर, अत्यंत गुंतागुंतीची
वाटणारी फलनेही अमूर्त अर्थाने संगणनक्षम असू शकतात.
उदाहरणार्थ, विश्वाचा कालावधी मर्यादित आहे असे समजा---काही
विश्वशास्त्रीय सिद्धांतांच्या भाकिताप्रमाणे फार दूरच्या भविष्यात
विश्व आकुंचन पावून एका बिंदूत सामावेल. मग त्या आरंभीच्या
क्षणापासून $f(n)$ परिभाषित असेल असे फक्त मर्यादित
(पण अविश्वसनीयरीत्या मोठ्या) संख्येने सेकंद असतील.
आणि केवळ मर्यादित संख्येच्या आदानांवर परिभाषित असलेले
कोणतेही फलन आंशिक संगणनक्षम असते: त्याची फलिते एका मोठ्या
सारणीमध्ये नोंदवता येतील किंवा ट्यूरिंग यंत्राच्या एका
अतिशय मोठ्या अवस्था-संक्रमण आलेखात सांकेतिक करता येतील.
यादीबाहेरच्या आदानांवर कार्यक्रम थांबणार नाही अशी व्यवस्था करता येते.


ज्या फलनांच्या किंमती होय/नाही या प्रश्नांची उत्तरे देतात,
अशा फलनांच्या विशेष प्रकारांत आपल्याला अनेकदा रस असतो.
उदाहरणार्थ, “$n$ ही अविभाज्य संख्या आहे का?” हा प्रश्न
पुढील फलनाशी संबंधित आहे:
\[\fn{isprime}(n) = \begin{cases}
  1 & \text{जर\/ $n$ अविभाज्य असेल तर}\\
  0 & \text{अन्यथा.}
  \end{cases}
\]
एखाद्या होय/नाही प्रश्नाशी संबंधित $1/0$-मूल्यी फलन
प्रभावी रीतीने संगणनक्षम असेल, तर त्या प्रश्नाचा
\emph{प्रभावी रीतीने निर्णय करता येतो} असे आपण म्हणतो.

प्रभावी रीतीने संगणित करता न येणारी फलने किंवा प्रभावी
रीतीने निर्णय करता न येणाऱ्या समस्या अस्तित्वात आहेत,
हे गणिती पद्धतीने सिद्ध करण्यासाठी संगणनाचे विशिष्ट
प्रतिमान निश्चित करणे आवश्यक आहे. त्यानंतर त्या
प्रतिमानाला काही फलने संगणित करता येत नाहीत किंवा
काही समस्यांचा निर्णय करता येत नाही, हे दाखवावे लागते.
उदाहरणार्थ, सर्वच फलने ट्यूरिंग यंत्रांनी संगणित करता
येत नाहीत आणि सर्वच समस्यांचा निर्णय ट्यूरिंग यंत्रांनी
करता येत नाही, हे आपण दाखवू शकतो. मग चर्च--ट्यूरिंग
प्रबंधाच्या आधारे केवळ ट्यूरिंग यंत्रे सर्व फलने संगणित
करण्याइतकी सामर्थ्यवान नाहीत असेच नव्हे, तर कोणतीही
प्रभावी प्रक्रियाही तसे करू शकत नाही, असा निष्कर्ष
काढता येतो.

असे निषेधात्मक निष्कर्ष सिद्ध करण्याची गुरुकिल्ली म्हणजे
ट्यूरिंग यंत्रांनाही संख्या देता येतात ही गोष्ट. याचा
सर्वांत सोपा मार्ग म्हणजे त्यांचे प्रगणन करणे: ट्यूरिंग
यंत्रे आणि त्यांचे कार्यक्रम लिहिण्याची एखादी निश्चित
पद्धत ठरवून त्यांची पद्धतशीर यादी करता येते. असे
करता येते हे लक्षात आल्यावर, केवळ संचसंख्येच्या
विचारांवरून ट्यूरिंग-असंगणनक्षम फलने अस्तित्वात आहेत
हे कळते: $\Nat$ पासून~$\Nat$ मध्ये जाणाऱ्या फलनांचा
संच (खरे तर, $\Nat$ पासून $\{0, 1\}$ मध्ये जाणाऱ्या
फलनांचा संचसुद्धा) अगणनीय म्हणजे अगणनीय
आहे; परंतु सर्व ट्यूरिंग यंत्रांचे प्रगणन करता येत
असल्यामुळे ट्यूरिंग-संगणनक्षम फलनांचा संच केवळ
गणनीय अनंत म्हणजे गणनीय अनंत आहे.

याशिवाय, संगणित करता येत नाहीत किंवा ज्यांचा निर्णय
करता येत नाही अशा अनुक्रमे \emph{विशिष्ट} फलनांची
आणि समस्यांचीही व्याख्या करता येते. त्यांपैकी
एक म्हणजे तथाकथित \emph{थांबण्याची समस्या.}
ट्यूरिंग यंत्रांच्या सूचना यादीत दिल्यास त्यांचे
मर्यादित वर्णन करता येते. ट्यूरिंग यंत्राचे असे
वर्णन, म्हणजेच ट्यूरिंग यंत्राचा कार्यक्रम, अर्थातच
दुसऱ्या ट्यूरिंग यंत्राला आदान म्हणून देता येते.
म्हणून इतर ट्यूरिंग यंत्रांबद्दलच्या प्रश्नांचा
निर्णय करणाऱ्या ट्यूरिंग यंत्रांचा विचार करता येतो.
विशेष रोचक प्रश्न असा आहे: “दिलेल्या ट्यूरिंग यंत्राला
आदान~$n$ दिल्यावर ते अखेरीस थांबते का?”
या प्रश्नाचा निर्णय करू शकणारे ट्यूरिंग यंत्र असते
तर बरे झाले असते: ट्यूरिंग यंत्रे अनंत चक्रात अडकत
नाहीत वगैरे तपासणारे गुणवत्ता-नियंत्रण यंत्र म्हणून
त्याचा विचार करा. ट्यूरिंग यांनी सिद्ध केलेली
रोचक गोष्ट म्हणजे असे यंत्र असूच शकत नाही.
ट्यूरिंग यंत्र~$M$ चे वर्णन आणि एखादी
संख्या~$n$ यांचे आदान दिल्यावर $M$ हे यंत्र
आदान $n$ वर थांबले असते की नाही यानुसार नेहमी
थांबून अनुक्रमे $1$ किंवा $0$ फलित देणारे एकही
ट्यूरिंग यंत्र असू शकत नाही.

अनिर्णेय समस्यांची विशिष्ट उदाहरणे मिळाल्यावर इतर
समस्याही अनिर्णेय आहेत हे दाखवण्यासाठी त्यांचा
उपयोग करता येतो. उदाहरणार्थ, “प्रथम-क्रम
सूत्र~$\varphi$ वैध आहे का?” हा एक प्रसिद्ध
अनिर्णेय प्रश्न आहे. आदान म्हणून प्रथम-क्रम
सूत्र~$\varphi$ दिल्यावर $\varphi$ वैध आहे की
नाही यानुसार $1$ किंवा $0$ फलित देत नेहमी
थांबेल असे कोणतेही ट्यूरिंग यंत्र नाही.
ऐतिहासिकदृष्ट्या, ही समस्या प्रभावी रीतीने
सोडवण्याची प्रक्रिया शोधण्याच्या प्रश्नाला
सरळ “निर्णय समस्या” म्हटले जात असे; म्हणून
निर्णय समस्या सोडवता येत नाही असे आपण म्हणतो.
ट्यूरिंग आणि चर्च यांनी हे फलित साधारणतः
एकाच काळात स्वतंत्रपणे सिद्ध केले; त्यामुळे
त्याला चर्च--ट्यूरिंग प्रमेय असेही म्हणतात.











\section{ट्यूरिंग यंत्रांचे प्रगणन}\label{tur:und:enu:sec}

\begin{explain}
सर्व ट्यूरिंग यंत्रांचा संच गणनीय म्हणजे गणनीय
आहे, असे आपण दाखवू शकतो. प्रत्येक ट्यूरिंग यंत्राचे
मर्यादित वर्णन करता येते, ही त्यामागची बाब आहे. अवस्थांचा
संच आणि फितीवरील चिन्हसंच हे दोन्ही सांत संच आहेत.
संक्रमण फलन हे $Q \times \Sigma$ पासून
$Q \times \Sigma \times \{\TMleft, \TMright, \TMstay\}$
मधील आंशिक फलन आहे. म्हणून ते परिभाषित असलेल्या
मर्यादित संख्येच्या आदान-जोड्यांवरील त्याची मूल्ये
यादीत दिल्यास त्याचेही वर्णन करता येते.

ही मांडणी एका मर्यादेपर्यंत खरी असली, तरी त्यात
एक सूक्ष्म फरक आहे. ट्यूरिंग यंत्राच्या व्याख्येत
अवस्थांच्या संचात आणि फितीच्या वर्णमालेत कोणते
घटक असावेत यावर बंधन घातलेले नाही.
म्हणून उदाहरणार्थ, प्रत्येक वास्तव संख्येला अवस्था
म्हणून वापरणारे ट्यूरिंग यंत्र तांत्रिकदृष्ट्या
असू शकते. मात्र ट्यूरिंग यंत्राचे \emph{वर्तन}
कोणत्या वस्तू अवस्था किंवा फितीवरील चिन्हे म्हणून
वापरल्या जातात यावर अवलंबून नसते.
\ref{tur:und:enu:fig:variants} मधील दोन ट्यूरिंग यंत्रे पाहा.
\begin{figure}
\begin{center}
\begin{tikzpicture}[->,>=stealth',shorten >=1pt,auto,node distance=2.8cm,
                    semithick]
  \tikzstyle{every state}=[fill=none,draw=black,text=black]

  \node[initial,state]         (A)              {$q_0$};
  \node[state]         (B) [right of=A] {$q_1$};

  \path (A) edge [bend left] node {\TMtrans{\TMstroke}{\TMstroke}{\TMright}} (B)
        (B) edge [loop above] node {\TMtrans{\TMblank}{\TMblank}{\TMright}} (B)
            edge [bend left] node {\TMtrans{\TMstroke}{\TMstroke}{\TMright}} (A);
\end{tikzpicture}\\
\begin{tikzpicture}[->,>=stealth',shorten >=1pt,auto,node distance=2.8cm,
  semithick]
\tikzstyle{every state}=[fill=none,draw=black,text=black]

\node[initial,state]         (A)              {$s$};
\node[state]         (B) [right of=A] {$h$};

\path (A) edge [bend left] node {\TMtrans{A}{A}{\TMright}} (B)
(B) edge [loop above] node {\TMtrans{\TMblank}{\TMblank}{\TMright}} (B)
edge [bend left] node {\TMtrans{A}{A}{\TMright}} (A);
\end{tikzpicture}
\end{center}
\caption{\emph{सम} यंत्राची रूपे}
\label{tur:und:enu:fig:variants}
\end{figure}
हे दोन आलेख दोन यंत्रे दर्शवतात: $M$ ची फितीची
वर्णमाला $\Sigma = \{\TMendtape,\TMblank,\TMstroke\}$
आणि अवस्थांचा संच $\{q_0,q_1\}$ आहे; तर $M'$ ची
वर्णमाला $\Sigma' = \{\TMendtape,\TMblank,A\}$ आणि
अवस्था $\{s,h\}$ आहेत. पण त्यांच्यातील सूचना
इतर बाबतीत सारख्याच आहेत: $n$ $\TMstroke$
चिन्हांची क्रमिका दिल्यावर $n$ सम असेल तेव्हा आणि केवळ तेव्हाच
$M$ थांबते; आणि $n$ $A$ चिन्हांची क्रमिका दिल्यावर
$n$ सम असेल तेव्हा आणि केवळ तेव्हाच $M'$ थांबते. आपण फक्त
$\TMstroke$ चे नाव~$A$, $q_0$ चे नाव~$s$
आणि $q_1$ चे नाव~$h$ असे बदलले आहे. हे
उदाहरण सर्वसाधारण करता येते: चिन्हे आणि अवस्था
अशा रीतीने पुनर्नामित केल्याने एक यंत्र दुसऱ्यात
रूपांतरित होत असेल, तर त्या ट्यूरिंग यंत्रांना
एकच मानता येते. खरे तर, ट्यूरिंग यंत्राची चिन्हे
आणि अवस्था या सरळ धन पूर्णांकच मानता येतात:
$\sigma_0$ ऐवजी~$1$, $\sigma_1$ ऐवजी~$2$
इत्यादी; $\TMendtape$ म्हणजे~$1$, $\TMblank$
म्हणजे~$2$ इत्यादी. अशा प्रकारे, \emph{सम}
यंत्र \ref{tur:und:enu:fig:standard-even} मध्ये दाखवलेल्या
यंत्रात रूपांतरित होते.
\begin{figure}
\[\begin{tikzpicture}[->,>=stealth',shorten >=1pt,auto,node distance=2.8cm,
  semithick]
\tikzstyle{every state}=[fill=none,draw=black,text=black]
\node[initial,state]         (A)              {$1$};
\node[state]         (B) [right of=A] {$2$};
\path (A) edge [bend left] node {\TMtrans{3}{3}{\TMright}} (B)
(B) edge [loop above] node {\TMtrans{2}{2}{\TMright}} (B)
edge [bend left] node {\TMtrans{3}{3}{\TMright}} (A);
\end{tikzpicture}
\]
\caption{एक मानक \emph{सम} यंत्र}
\label{tur:und:enu:fig:standard-even}
\end{figure}
ज्या ट्यूरिंग यंत्राच्या अवस्था आणि चिन्हे धन पूर्णांक
असतात त्याला \emph{मानक} यंत्र म्हणता येईल; यापुढे
केवळ मानक यंत्रांचाच विचार करू.\footnote{“मानक
यंत्र” ही संज्ञा स्वतः मानक संज्ञा नाही.}

ट्यूरिंग यंत्रांचा संच गणनीय म्हणजे गणनीय
आहे, हे दाखवण्याचे आपले उद्दिष्ट होते. वरील विचार
लक्षात घेतल्यास मानक ट्यूरिंग यंत्रांचा संच
गणनीय म्हणजे गणनीय
आहे हे दाखवणे पुरेसे आहे. समजा, एक मानक ट्यूरिंग
यंत्र $M = \tuple{Q, \Sigma, q_0, \delta}$ दिले
आहे. धन पूर्णांकांच्या एका सांत क्रमिकेत त्याचे
वर्णन कसे करावे? प्रथम अवस्थांची संख्या, मग
अवस्था, चिन्हांची संख्या, चिन्हे आणि आरंभीची
अवस्था अशी यादी करू. (लक्षात घ्या, $M$ मानक
यंत्र असल्यामुळे ही सर्व धन पूर्णांक आहेत.)
मग~$\delta$ चे काय? संभाव्य आदाने, म्हणजे
$\tuple{q,\sigma}$ या जोड्या, यांचा संच सांत आहे;
कारण $Q$ आणि~$\Sigma$ सांत आहेत. म्हणून~$\delta$
मधील माहिती म्हणजे ज्या सर्व बाबतीत
$\delta(q,\sigma) = \tuple{q', \sigma', D}$ आहे,
अशा $\tuple{q, \sigma, q', \sigma', d}$
या सर्व $5$-घटकी क्रमितांची सांत यादी होय.
येथे $d$ ही दिशा~$D$ सांकेतिक
करणारी संख्या आहे (उदाहरणार्थ,
~$\TMleft$ साठी $1$, ~$\TMright$ साठी $2$
आणि ~$\TMstay$ साठी $3$).

अशा प्रकारे प्रत्येक मानक ट्यूरिंग यंत्राचे
धन पूर्णांकांच्या सांत यादीने, म्हणजे
$s_M \in (\PosInt)^*$ या क्रमिकेने, वर्णन
करता येते. उदाहरणार्थ, मानक \emph{सम}
यंत्र पुढील क्रमिकेने सांकेतिक केले जाते:
\[2, \underbrace{1, 2}_Q, 3, \overbrace{1, 2, 3}^\Sigma, 1, \underbrace{1, 3, 2, 3, 2}_{\delta(1,3) = \tuple{2,3,R}},
\overbrace{2, 2, 2, 2, 2}^{\delta(2,2) = \tuple{2,2,R}},
\underbrace{2, 3, 1, 3, 2}_{\delta(2,3) = \tuple{1,3,R}}.
\]
\end{explain}

\begin{thm}
$\Nat$ पासून~$\Nat$ मध्ये जाणारी अशी फलने
आहेत जी ट्यूरिंग-संगणनक्षम नाहीत.
\end{thm}

\begin{proof}
धन पूर्णांकांच्या सांत क्रमिकांचा संच~$(\PosInt)^*$
गणनीय म्हणजे गणनीय आहे, हे आपल्याला माहीत आहे
(\ref{sfr:siz:zigzag:prob:posint-star}). मानक ट्यूरिंग
यंत्रांच्या वर्णनांचा संच हा~$(\PosInt)^*$ चा उपसंच
असल्याने तोही गणनीय आहे. $\Nat$ पासून~$\Nat$
मध्ये जाणारे प्रत्येक ट्यूरिंग-संगणनक्षम फलन कोणत्या
ना कोणत्या (खरे तर अनेक) ट्यूरिंग यंत्रांनी संगणित
केले जाते. त्याच्या अवस्था आणि चिन्हांना धन
पूर्णांकांची नावे देऊन (विशेषतः $\TMendtape$ ला~$1$,
$\TMblank$ ला~$2$ आणि $\TMstroke$ ला~$3$)
प्रत्येक ट्यूरिंग-संगणनक्षम फलन कोणत्या ना कोणत्या
मानक ट्यूरिंग यंत्रानेही संगणित केले जाते, हे
दिसते. म्हणून $\Nat$ पासून~$\Nat$ मध्ये जाणाऱ्या
सर्व ट्यूरिंग-संगणनक्षम फलनांचा संचही गणनीय आहे.

दुसरीकडे, $\Nat$ पासून~$\Nat$ मध्ये जाणाऱ्या
सर्व फलनांचा संच गणनीय म्हणजे गणनीय नाही
(\ref{sfr:siz:red:prob:nat-nat}). जर प्रत्येक
फलन एखाद्या ट्यूरिंग यंत्राने संगणित केले
जात असेल, तर ती फलने संगणित करणाऱ्या यंत्रांची
सर्व वर्णने यादीत दिल्याने सर्व फलनांचे प्रगणन
करता आले असते. त्यामुळे काही फलने
ट्यूरिंग-संगणनक्षम नाहीत.
\end{proof}

\begin{prob}
  अवस्था आणि वर्णमालेतील चिन्हांची स्वतंत्र
  स्पष्ट यादी न देता ट्यूरिंग यंत्रांचे वर्णन
  करण्याचा मार्ग तुम्हाला सुचतो का? “मानक”
  यंत्राची स्वतःची संकल्पना तुम्ही ठरवू शकता;
  पण तुमच्या नव्या अर्थाने प्रत्येक ट्यूरिंग
  यंत्राचे संगणन एखाद्या “मानक” यंत्राने
  कसे करता येईल, हेही स्पष्ट करा.
\end{prob}











\section{सार्वत्रिक ट्यूरिंग यंत्रे}\label{tur:und:uni:sec}

\ref{tur:und:enu:sec} मध्ये प्रत्येक ट्यूरिंग यंत्राचे
पूर्णांकांच्या सांत क्रमिकेने वर्णन कसे करता येते,
याची चर्चा केली. ही क्रमिका त्या यंत्राच्या अवस्था,
वर्णमाला, आरंभीची अवस्था आणि सूचना सांकेतिक करते.
अशा सर्व वर्णनांचा संच गणनीय म्हणजे गणनीय
आहे, हेही आपण दाखवले. अशा वर्णनांचा संच
गणनीय अनंत म्हणजे गणनीय अनंत असल्यामुळे
~$\Nat$ पासून या वर्णनांवर जाणारे
आच्छादक म्हणजे आच्छादक फलन आहे.
असे आच्छादक फलन उदाहरणार्थ कांतोरची
नागमोडी पद्धत वापरून मिळवता येते. त्यामुळे
ट्यूरिंग यंत्रांच्या (वर्णनांचे) प्रगणन करता येते.
असे एक प्रभावी प्रगणन शून्य निर्देशांकापासून निश्चित केल्यावर $0$-वे, $1$-वे,
$2$-रे, \dots, $e$-वे ट्यूरिंग यंत्र असे
बोलता येते. या संख्यांना \emph{निर्देशांक} म्हणतात.


\begin{defn}
$M$ हे (आपण निश्चित केलेल्या प्रगणनातील)
$e$-वे ट्यूरिंग यंत्र असेल, तर $e$ हा~$M$ चा
\emph{निर्देशांक} आहे असे म्हणतो. $e$-व्या
ट्यूरिंग यंत्रासाठी $M_e$ असे लिहितो.
\end{defn}

एका यंत्राला एकाहून अधिक निर्देशांक असू शकतात.
उदाहरणार्थ, $M$~च्या सूचनांची यादी कोणत्या
क्रमाने करतो यावरून त्याची दोन वर्णने वेगळी
असू शकतात; अशा वेगळ्या वर्णनांचे निर्देशांकही
वेगवेगळे असतात.

महत्त्वाची बाब अशी की ट्यूरिंग यंत्रांच्या वर्णनांचे
प्रगणन अशा रीतीने करता येते की निर्देशांकावरून
यंत्र~$M$ चे वर्णन प्रभावी रीतीने संगणित करता
येईल आणि यंत्र~$M$ च्या वर्णनावरून त्याचा
एक निर्देशांकही प्रभावी रीतीने संगणित करता
येईल. मग चर्च--ट्यूरिंग प्रबंधानुसार, निर्देशांक~$e$
असलेल्या ट्यूरिंग यंत्राचे वर्णन मिळवून ते
आपल्या फितीवर फलित म्हणून लिहिणारे ट्यूरिंग
यंत्र असू शकते. हे वर्णन म्हणजे
~$\TMstroke$ चिन्हांच्या खंडांची क्रमिका असेल
(हे खंड $M_e$ चे वर्णन करणाऱ्या क्रमिकेतील
धन पूर्णांक दर्शवतील).

आता असा प्रश्न स्वाभाविकपणे पडतो: ट्यूरिंग
यंत्रांच्या निर्देशांकांवरील कोणती फलने स्वतः
ट्यूरिंग यंत्रांनी संगणित करता येतात? ट्यूरिंग
यंत्रांच्या निर्देशांकांचे कोणते गुणधर्म ट्यूरिंग
यंत्रांनी निर्णीत करता येतात? उदाहरण पाहा:
निर्देशांक~$e$ वरून निर्देशांक~$e$ असलेल्या
ट्यूरिंग यंत्रातील अवस्थांची संख्या देणारे
फलन ट्यूरिंग यंत्राने संगणित करता येते.
असे यंत्र पुढीलप्रमाणे काम करेल: सुरुवातीला
त्याच्या फितीवर $e$~$\TMstroke$ चिन्हांचा
एकच खंड असेल. प्रथम ते $e$ वरून त्याचे
वर्णन उलगडेल. आता ते वर्णन फितीवरील
~$\TMstroke$ चिन्हांच्या खंडांच्या क्रमिकेने
दर्शवलेले असेल. या क्रमिकेतील पहिला
घटक म्हणजे अवस्थांची संख्या आहे.
म्हणून पहिला $\TMstroke$ खंड सोडून इतर
सर्व पुसून टाकायचे आणि मग थांबायचे, एवढेच
उरते.

पुढील फलित विशेष उल्लेखनीय आहे:

\begin{thm}\label{tur:und:uni:thm:universal-tm} असे एक \emph{सार्वत्रिक
  ट्यूरिंग यंत्र}~$U$ आहे की आदान $\tuple{e,n}$ वर
  सुरू केल्यास ते
  \begin{enumerate}
    \item $M_e$ हे आदान~$n$ वर थांबते तेव्हा आणि केवळ तेव्हाच थांबते, आणि
    \item $M_e$ हे फलित $m$ देऊन थांबले, तर~$U$ देखील
    त्याच फलितासह थांबते.
  \end{enumerate}
  म्हणून $U$ हे $f\colon \Nat \times \Nat \pto \Nat$
  फलन संगणित करते; $M_e$ हे आदान~$n$ वर सुरू
  होऊन फलित~$m$ देत थांबले तर $f(e,n) = m$,
  आणि अन्यथा ते अपरिभाषित असते.
\end{thm}

\begin{proof}
  $U$ चा संपूर्ण सूचनासंच येथे प्रत्यक्ष लिहून काढणे अव्यवहार्य आहे, कारण
  या रचनेतील यंत्र अत्यंत गुंतागुंतीचे आहे. पण ते कसे
  काम करते याची रूपरेषा देता येते आणि मग
  चर्च--ट्यूरिंग प्रबंधाचा आधार घेता येतो.
  सुरुवातीला $U$ च्या फितीवर $e$ $\TMstroke$
  चिन्हांचा खंड आणि त्यानंतर $n$~$\TMstroke$
  चिन्हांचा खंड असतो. ते प्रथम निर्देशांक~$e$
  आदान~$n$ च्या उजवीकडे “उलगडते.” यामुळे
  यंत्र~$M_e$ च्या सूचनांचे वर्णन करणाऱ्या
  संख्यांची यादी मिळते (म्हणजे~$\TMblank$
  चिन्हांनी वेगळे केलेले $\TMstroke$ चिन्हांचे
  खंड). त्यानंतर $U$ हे $M_e$ च्या वर्णनाच्या
  उजवीकडे $M_e$ च्या आरंभीच्या अवस्थेची
  संख्या आणि संख्या~$1$ लिहिते. (पुन्हा,
  या संख्या $\TMstroke$ चिन्हांच्या खंडांनी
  एकचिन्ही रूपात दर्शवलेल्या आहेत.) पुढे,
  आदानातील $n$~$\TMstroke$ चिन्हांचा खंड
  उजवीकडे नकलते---पण प्रत्येक $\TMstroke$
  चिन्हाच्या जागी तीन $\TMstroke$ चिन्हांचा
  खंड लिहिते ($\TMstroke$ या चिन्हाची संख्या~$3$
  आहे, तर~$\TMendtape$ ची संख्या $1$ आणि
  ~$\TMblank$ ची संख्या $2$ आहे, हे आठवा).
  अशा खंडांच्या क्रमिकेच्या डाव्या टोकाला
  ($U$ च्या फितीवर ते~$\TMblank$ चिन्हांनी
  वेगळे केलेले आहेत) ते एकच~$\TMstroke$
  चिन्ह लिहिते; हाच~$\TMendtape$ चा संकेतांक आहे.
  आदान रिकामे असेल, तर त्यानंतरच्या पहिल्या रिकाम्या घराचाही
  सांकेतिक खंड लिहिते, जेणेकरून प्रारंभिक शीर्षासाठी घर उपलब्ध असेल.

  आता $U$ च्या फितीवर निर्देशांक~$e$, संख्या~$n$,
  आरंभीच्या अवस्थेचा संकेतांक (“सध्याची अवस्था”),
  आरंभीच्या शीर्षस्थानाची संख्या~$1$
  (“सध्याचे शीर्षस्थान”) आणि अनुकरणातील
  “फिती”वरील आरंभीचा मजकूर नोंदलेला असतो
  (म्हणजे~$M_e$ च्या “फिती”वरील “चिन्हांचे”
  संकेतांक दर्शवणाऱ्या~$\TMstroke$ चिन्हांच्या
  खंडांची, ~$\TMblank$ चिन्हांनी वेगळे केलेली
  क्रमिका).

  आदान~$n$ वर सुरू केलेले $M_e$ जे काम
  करेल त्याचे अनुकरण ते पुढीलप्रमाणे करते:
  \begin{enumerate}
    \item “सध्याच्या शीर्षस्थानाची” संख्या~$k$ शोधणे
    (सुरुवातीला ती~$1$ असते),
    \item डाव्या टोकाचे चिन्ह दर्शवणारा खंड शून्य क्रमांकाचा धरून
    “फितीवरील” $k$-व्या खंडाकडे जाऊन
    तेथील “चिन्ह” कोणते आहे ते पाहणे,
    \item\label{tur:und:uni:find-inst}
    सध्याची “अवस्था” आणि “चिन्ह” यांना जुळणारी
    सूचना शोधणे,
    \item “फितीवरील” $k$-व्या खंडाकडे
    परत जाऊन तेथील “चिन्हा”ऐवजी $M_e$ लिहील
    त्या चिन्हाचा संकेतांक लिहिणे,
    \item सध्याची “अवस्था” नोंदवलेली आहे
    तेथे शीर्ष नेऊन त्या संख्येऐवजी नव्या
    अवस्थेची संख्या लिहिणे,
    \item “फितीवरील स्थान” नोंदवलेले आहे
    तेथे जाऊन एक~$\TMstroke$ पुसणे किंवा
    एक~$\TMstroke$ जोडणे (सूचना अनुक्रमे
    डावीकडे किंवा उजवीकडे जायला सांगत असेल तर). सध्याचे
    शीर्षस्थान शून्य असेल आणि डावीकडे जाण्याची सूचना असेल,
    तर ते स्थान बदलत नाही. उजवीकडे सरकल्यावर सांकेतिक फितीच्या
    सध्याच्या शेवटापलीकडे घर लागल्यास नवीन रिकाम्या घराचा खंड जोडते.
    \item पुन्हा हेच करणे.\footnote{येथे आपण काही सूक्ष्म
    अडचणींचा तपशील टाळत आहोत. उदाहरणार्थ,
    “सध्याचे शीर्षस्थान” मोजणाऱ्या संख्येत
    वाढ करताना $U$ ला अतिरिक्त जागा लागू शकते;
    अशा वेळी त्याला संपूर्ण “फीत” उजवीकडे
    हलवावी लागेल.}
  \end{enumerate}
  आदान~$n$ वर सुरू केलेले $M_e$ कधीच थांबत
  नसेल, तर $U$ देखील कधीच थांबत नाही;
  म्हणून त्याचे फलित अपरिभाषित असते.

  पायरी~\ref{tur:und:uni:find-inst} मध्ये $M_e$ च्या
  वर्णनात सध्याच्या “अवस्था”/“चिन्ह” जोडीला
  जुळणारी सूचना नसेल, तर $M_e$ थांबले असते.
  तसे झाल्यास $U$ प्रथम अनुकरणातील फितीवर प्रदानाची व्याख्या
  लागू होते का ते तपासते: सुरुवातीला मूळ डाव्या टोकाचे चिन्ह,
  त्यानंतर सलग रेघा आणि त्यानंतर फक्त रिकामी घरे असली पाहिजेत.
  ही अट पूर्ण असल्यास आपल्या फितीवरील “फिती”च्या
  डावीकडचा भाग पुसते. $M_e$ च्या फितीवरील
  एक~$\TMstroke$ दर्शवणाऱ्या प्रत्येक तीन
  ~$\TMstroke$ चिन्हांच्या खंडासाठी ते
  स्वतःच्या फितीच्या डाव्या टोकाच्या चिन्हाच्या लगेच उजवीकडे एक
  $\TMstroke$ लिहिते आणि त्याच वेळी क्रमाने
  “फीत” पुसत जाते. हे संपल्यावर $U$ च्या
  फितीवर लांबी~$m$ असलेला $\TMstroke$
  चिन्हांचा एकच खंड उरतो. रेघा नसल्यास प्रदान रिकामे असते.

  वरची पूर्ण प्रदान-अट मोडली असेल—उदाहरणार्थ, मूळ प्रदानातील
  रेघेसाठी तीन~$\TMstroke$ चिन्हांचा खंड नसणे, मध्ये रिकामे घर
  असणे किंवा आरंभीचे डाव्या टोकाचे चिन्ह नसणे—तर $U$
  एक रिकामे चिन्ह आणि त्यानंतर एक रेघ असा मुद्दाम अमान्य
  प्रदानाचा आकार लिहून थांबते. या बाबतीत $U$ च्या फितीवर
  $\TMstroke$ चिन्हांचा एकच खंड नसल्याने
  त्याचे फलित नैसर्गिक संख्या नसते; म्हणजे
  या बाबतीत $f(e,n)$ अपरिभाषित असते. आरंभीच्या डाव्या टोकाच्या
  चिन्हाचा सांकेतिक खंड आणि अखेरीच्या रिकाम्या घरांचे खंड
  हे वैध प्रदानात असू शकतात; ते रेघांचे खंड नाहीत म्हणून नाकारायचे नाहीत.

\end{proof}











\section{थांबण्याची समस्या}\label{tur:und:hal:sec}

\begin{explain}
ट्यूरिंग यंत्रांच्या वर्णनांचे एक प्रगणन आपण निश्चित
केले आहे असे समजा. त्यामुळे प्रत्येक ट्यूरिंग
यंत्राला एक \emph{निर्देशांक} मिळतो: ट्यूरिंग
यंत्रांच्या वर्णनांच्या $M_0$, $M_1$, $M_2$, $M_3$,
\dots{} या प्रगणनातील त्याचे स्थान.


ट्यूरिंग-संगणनक्षम नसलेली फलने असलीच पाहिजेत,
हे आपल्याला माहीत आहे: ट्यूरिंग यंत्रांच्या
वर्णनांचा---आणि म्हणून यंत्रांचा---संच
गणनीय म्हणजे गणनीय आहे; पण
$\Nat$ पासून $\Nat$ मध्ये जाणाऱ्या सर्व
फलनांचा संच तसा नाही. मात्र संगणनक्षम
नसलेल्या विशिष्ट फलनांची उदाहरणेही देता
येतात. त्यांपैकी एक म्हणजे थांबण्याचे फलन.
\end{explain}

\begin{defn}[थांबण्याचे फलन]
 \emph{थांबण्याचे फलन}~$h$ असे परिभाषित करतात:
\[h(e,n) =
\begin{cases}
  \text{0} & \text{जर यंत्र~$M_e$ आदान $n$ वर थांबत नसेल तर} \\
  \text{1} & \text{जर यंत्र~$M_e$ आदान $n$ वर थांबत असेल तर}
\end{cases}
\]
\end{defn}

\begin{defn}[थांबण्याची समस्या]
\emph{थांबण्याची समस्या} म्हणजे (कोणत्याही $e$, $n$
साठी) ट्यूरिंग यंत्र~$M_e$ हे $n$ रेघांच्या
आदानावर थांबते का हे ठरवण्याची समस्या.
\end{defn}

\begin{explain}
~$h$ हे ट्यूरिंग-संगणनक्षम नाही, हे आपण
त्याच्याशी संबंधित~$s$ हे फलन ट्यूरिंग-संगणनक्षम
नाही असे दाखवून सिद्ध करू. हा पुरावा कोणत्याही
ट्यूरिंग यंत्राचे संगणन शिस्तबद्ध ट्यूरिंग
यंत्राने करता येते (\ref{tur:mac:dis:sec}),
आणि दोन ट्यूरिंग यंत्रे जोडून एक यंत्र
बनवता येते (\ref{tur:mac:cmb:sec}),
या दोन गोष्टींवर आधारलेला आहे.
\end{explain}

\begin{defn} फलन~$s$ पुढीलप्रमाणे परिभाषित आहे:
\[s(e) =
\begin{cases}
  \text{0} & \text{जर यंत्र~$M_e$ आदान $e$ वर थांबत नसेल तर} \\
  \text{1} & \text{जर यंत्र~$M_e$ आदान $e$ वर थांबत असेल तर}
\end{cases}
\]
\end{defn}

\begin{lem}
फलन~$s$ ट्यूरिंग-संगणनक्षम नाही.
\end{lem}

\begin{proof}
विरोधासाठी फलन~$s$ ट्यूरिंग-संगणनक्षम आहे असे
समजा. मग~$s$ संगणित करणारे ट्यूरिंग यंत्र~$S$
असेल. व्यापकतेला बाधा न आणता, $S$ शिस्तबद्ध
आहे असे मानता येते. विशेषतः, ते थांबते
तेव्हा पहिल्या घराकडे पाहत असते.

या यंत्राला दुसरे यंत्र~$J$ जोडता येते.
त्याला आदान~$0$ दिल्यास ते थांबते
(म्हणजे आरंभीच्या अवस्थेत, फितीच्या टोकाच्या
चिन्हाच्या उजवीकडील घराकडे पाहताना त्याला
$\TMblank$ दिसल्यास); अन्यथा ते उजवीकडे
चालतच राहते आणि कधीच थांबत नाही.
$S$ ला~$J$ जोडून तयार केलेले
$S \concat J$ हे ट्यूरिंग यंत्र आहे;
म्हणून ते कोणत्यातरी~$e$ साठी $M_e$
आहे (म्हणजे प्रगणनात कुठेतरी येते).
$M_e$ हे $e$ $\TMstroke$ चिन्हांच्या
आदानावर सुरू करा. दोनच शक्यता आहेत:
$M_e$ थांबेल किंवा थांबणार नाही.
\begin{enumerate}
\item $M_e$ हे $e$ $\TMstroke$ चिन्हांच्या
  आदानावर थांबते असे समजा. मग $s(e)
  = 1$. म्हणून~$e$ वर सुरू केलेले $S$
  थांबते आणि फितीवर फलित म्हणून एकच
  $\TMstroke$ ठेवते. मग $J$ फितीवरील
  $\TMstroke$ चिन्हापासून सुरू होते. त्या
  बाबतीत $J$ थांबत नाही. पण $M_e$ हे
  $S \concat J$ यंत्रच आहे; म्हणून $S$
  नंतर~$J$ जे करेल तेच त्याने करायला
  हवे (म्हणजे या बाबतीत उजवीकडे चालत
  राहून कधीच न थांबणे). म्हणून $M_e$
  हे $e$ $\TMstroke$ चिन्हांच्या आदानावर
  थांबू शकत नाही.

\item आता $M_e$ हे $e$ $\TMstroke$
  चिन्हांच्या आदानावर थांबत नाही असे
  समजा. मग $s(e) = 0$ आणि आदान~$e$
  वर सुरू केलेले $S$ रिकामी फीत ठेवून
  थांबते. रिकाम्या फितीवर सुरू केलेले
  $J$ लगेच थांबते. पुन्हा, $M_e$ हे
  $S$ नंतर~$J$ जे करेल तेच करते;
  म्हणून $M_e$ हे $e$ $\TMstroke$
  चिन्हांच्या आदानावर थांबलेच पाहिजे.
\end{enumerate}
दोन्ही बाबतीत आपल्या गृहीतकाला विरोध
मिळतो. त्यामुळे असे ट्यूरिंग यंत्र~$S$
असू शकत नाही: $s$ ट्यूरिंग-संगणनक्षम नाही.
\end{proof}

\begin{thm}[थांबण्याच्या समस्येची असमाधेयता]
\label{tur:und:hal:thm:halting-problem} थांबण्याची समस्या
सोडवता येत नाही; म्हणजे फलन~$h$
ट्यूरिंग-संगणनक्षम नाही.
\end{thm}

\begin{proof}
$h$ हे ट्यूरिंग-संगणनक्षम आहे, आणि समजा
ट्यूरिंग यंत्र~$H$ ते संगणित करते.
मग~$H$ वापरून~$s$ संगणित करणारे ट्यूरिंग
यंत्र बनवता येईल: प्रथम आदानाची एक प्रत
बनवा (दोन्ही प्रतींमध्ये~$\TMblank$ चिन्ह
ठेवा). मग पुन्हा आरंभी जा आणि~$H$
चालवा. पहिले काम करणारे यंत्र स्पष्टपणे
बनवता येते (\ref{tur:mac:dis:prob:copier}
पाहा); आणि $H$ अस्तित्वात असेल, तर
ते अशा प्रतिकृती तयार करणाऱ्या यंत्राला
जोडून $M_e$ हे आदान~$e$ वर थांबते
का हे ठरवणारे, म्हणजे~$s$ संगणित
करणारे, नवे यंत्र मिळेल. पण असे यंत्र
असू शकत नाही, हे आपण आधीच दाखवले
आहे. म्हणून~$h$ देखील ट्यूरिंग-संगणनक्षम
नाही.
\end{proof}

\begin{prob}
तीन-रेघांची थांबण्याची (3-Halt) समस्या
म्हणजे अन्यथा रिकाम्या असलेल्या फितीवर
तीन $\TMstroke$ चिन्हांचे आदान दिल्यावर
स्वैरपणे निवडलेले ट्यूरिंग यंत्र थांबते
की नाही हे ठरवणारी निर्णय प्रक्रिया
देण्याची समस्या. 3-Halt समस्या
सोडवता येत नाही हे सिद्ध करा.
\end{prob}

\begin{prob}
ट्यूरिंग यंत्र आणि आदान यांच्या
$M_e$ आणि~$n$ या जोड्यांसाठी,
$e \neq n$ असताना थांबण्याची समस्या
सोडवता येत असेल, तर $e = n$
असलेल्या बाबतीतही ती सोडवता
येते, हे दाखवा.
\end{prob}

\begin{prob}
आदान~$e$ ही संख्या असेल आणि ती
सर्व ट्यूरिंग यंत्रांच्या प्रगणनाद्वारे
यंत्र~$M_e$ ओळखून देत असेल,
तर थांबण्याची समस्या सोडवता
येत नाही, हे आपण सिद्ध केले.
त्याऐवजी \ref{tur:und:enu:sec}
मधील ट्यूरिंग यंत्रांची वर्णनेच
थेट आदान म्हणून दिली तर काय?
निर्देशांकांऐवजी ट्यूरिंग यंत्रांची
वर्णने आदान म्हणून घेऊन थांबण्याच्या
समस्येचा निर्णय करणारे ट्यूरिंग
यंत्र असू शकते का? का किंवा
का नाही, ते स्पष्ट करा.
\end{prob}

\begin{prob} पुढीलप्रमाणे परिभाषित
  \emph{आंशिक} फलन~$s'$ हे
  ट्यूरिंग-संगणनक्षम \emph{आहे},
  हे दाखवा:
  \[  s'(e) =
  \begin{cases}
    \text{1} & \text{जर यंत्र~$M_e$ आदान $e$ वर थांबत असेल तर}\\
    \text{अपरिभाषित} & \text{जर यंत्र~$M_e$ आदान $e$ वर थांबत नसेल तर}
  \end{cases}
  \]
\end{prob}











\section{निर्णय समस्या}\label{tur:und:dec:sec}

दिलेल्या वाक्याची वैधता ठरवण्याची
प्रभावी पद्धत असेल तेव्हा आणि केवळ तेव्हाच
प्रथम-क्रम तर्कशास्त्र \emph{निर्णेय} आहे,
असे आपण म्हणतो. प्रत्यक्षात अशी पद्धत नाही:
प्रथम-क्रम वाक्यांच्या वैधतेचा निर्णय
करण्याची समस्या सोडवता येत नाही.

हा महत्त्वाचा निषेधात्मक निष्कर्ष स्थापित
करण्यासाठी, निर्णय समस्या ट्यूरिंग यंत्राने
सोडवता येत नाही हे आपण सिद्ध करतो.
म्हणजे, फितीवर प्रथम-क्रम वाक्य असताना सुरू केले की ते वाक्य वैध
आहे की नाही यानुसार अखेरीस थांबून
$1$ किंवा~$0$ फलित देईल, असे कोणतेही
ट्यूरिंग यंत्र नाही हे दाखवतो.
चर्च--ट्यूरिंग प्रबंधानुसार, संगणनक्षम
असलेले प्रत्येक फलन ट्यूरिंग-संगणनक्षम
असते. म्हणून हे “वैधतेचे फलन”
मुळात प्रभावी रीतीने संगणित करता
येत असेल, तर ते ट्यूरिंग-संगणनक्षमही
असेल. ते ट्यूरिंग-संगणनक्षम नसेल,
तर प्रभावी रीतीनेही संगणित करता
येणार नाही.

निर्णय समस्या सोडवता येत नाही हे
सिद्ध करण्यासाठी आपण थांबण्याच्या
समस्येचे तिच्याकडे न्यूनीकरण करू.
याचा अर्थ पुढीलप्रमाणे. ट्यूरिंग
यंत्राचे वर्णन~$e$ या संकेतांकाने होत असेल,
तर ते यंत्र आदान~$w$ वर थांबते
तेव्हा~$1$ आणि अन्यथा~$0$ देणारे
फलन~$h(e,w)$ ट्यूरिंग-संगणनक्षम
नाही: आधी सिद्ध केलेली एकचिन्ही नैसर्गिक-संख्या आदानांची
थांबण्याची समस्या हे याचे विशेष प्रकरण आहे.
प्रथम-क्रम वाक्यांची वैधता ठरवणारे
ट्यूरिंग यंत्र असेल, तर~$h$ संगणित
करणारे ट्यूरिंग यंत्रही असेल, हे
आपण दाखवू. $h$ ट्यूरिंग यंत्राने
संगणित करता येत नसल्यामुळे वैधतेचा
निर्णय करणारे ट्यूरिंग यंत्रही असू
शकत नाही.

या युक्तीतील पहिली पायरी अशी आहे:
प्रत्येक आदान~$w$ आणि ट्यूरिंग
यंत्र~$M$ साठी, $M$ चा सूचनासंच
आणि आदान~$w$ दर्शवणारे
वाक्य $\mathsf{T}(M, w)$,
तसेच “$M$ अखेरीस थांबते”
हे व्यक्त करणारे वाक्य~$\mathsf{E}(M, w)$ प्रभावी रीतीने
लिहिता येतात, असे दाखवायचे;
आणि ती अशी असतील की
\begin{quote}
  $\Entails \mathsf{T}(M, w) \lif \mathsf{E}(M,w)$ हे खरे असते तेव्हा आणि केवळ तेव्हाच, जेव्हा $M$ हे आदान~$w$ वर थांबते.
\end{quote}
आपल्या पुराव्याचा मुख्य भाग ही
वाक्ये $\mathsf{T}(M, w)$ आणि~$\mathsf{E}(M, w)$
यांचे वर्णन करण्यात आणि
$\mathsf{T}(M, w) \lif \mathsf{E}(M, w)$ हे
आदान~$w$ वर~$M$ थांबते तेव्हा आणि केवळ तेव्हाच
वैध असते हे पडताळण्यात जाईल.











\section{ट्यूरिंग यंत्रांचे निरूपण}\label{tur:und:rep:sec}

\begin{explain}
ट्यूरिंग यंत्रे आणि त्यांचे वर्तन प्रथम-क्रम तर्कशास्त्रातील
वाक्याने निरूपित करण्यासाठी योग्य भाषा
परिभाषित करावी लागते. त्या भाषेचे दोन भाग आहेत:
यंत्राच्या स्थितिवर्णनांचे वर्णन करणारी विधेये
म्हणजे विधेये, आणि संगणनाच्या पायऱ्या (“क्षण”)
तसेच फितीवरील स्थाने यांना संख्या देण्यासाठीचे
पदबंध.

दोन प्रकारची विधेये म्हणजे विधेये
आपण घेऊ; दोन्ही द्विस्थानी आहेत: प्रत्येक
अवस्था~$q$ साठी विधेय~$\Obj Q_q$,
आणि फितीवरील प्रत्येक चिन्ह~$\sigma$
साठी विधेय~$\Obj S_\sigma$.
पहिल्या प्रकाराने~$M$ ची अवस्था आणि
त्याच्या फितीवरील शीर्षाचे स्थान वर्णन
करता येते; दुसऱ्या प्रकाराने फितीवरील
मजकूर वर्णन करता येतो.

फितीवरील शीर्षाचे स्थान आणि केलेल्या
पायऱ्यांची संख्या व्यक्त करण्यासाठी
संख्या व्यक्त करण्याचा मार्ग हवा.
यासाठी स्थिरांक~$\Obj 0$
आणि उत्तरवर्ती फलन~$\prime$ हे
$1$-स्थानी फलन वापरतो. संकेताप्रमाणे
ते त्याच्या आदानाच्या \emph{नंतर}
लिहितो (आणि कंस गाळतो).

प्रत्येक संख्या $n$ साठी तिचे
~$\Lang L_M$ मधील निरूपण करणारे
मानक पद~$\num{n}$ असते; त्याला
~$n$ चा \emph{संख्यांक} म्हणतात.
$\num{0}$ म्हणजे $\Obj 0$,
$\num{1}$ म्हणजे $\Obj 0'$,
$\num{2}$ म्हणजे $\Obj 0''$,
इत्यादी. अधिक औपचारिकरीत्या:
\begin{align*}
\num{0} & = \Obj 0 \\
\num{n+1} &= \num{n}'
\end{align*}
$\num{0}$ हे पद, म्हणजे $\Obj 0$,
फितीवरील सर्वांत डावीकडील घराचे स्थान
आणि संगणनाची पहिली पायरी सुरू होण्यापूर्वीचा
क्षण (प्रारंभिक स्थितिवर्णन) या दोन्हींना
नाव देते. $\num{1}$ हे पद, म्हणजे
$\Obj 0'$, सर्वांत डावीकडील घराच्या
उजवीकडील घराला आणि संगणनाच्या
पहिल्या पायरीनंतरच्या क्षणाला नाव
देते; पुढेही याचप्रमाणे.

फितीवरील स्थानांचा क्रम (“याच्या डावीकडे”)
आणि संगणनाच्या पायऱ्यांचा क्रम (“याच्या
आधी”) हे दोन्ही व्यक्त करण्यासाठी
विधेय~$<$ सुद्धा घेऊ.

भाषा ठरल्यावर $\mathsf{T}(M,
w)$ ची “स्वयंसिद्धके” यादीत देऊ.
म्हणजे आदान~$w$ वर चालवलेल्या~$M$
चे वर्तन एकत्रितपणे वर्णन करणारी
वाक्ये. त्यांत
$\Obj 0$, $\prime$ आणि $<$ यांवर
अटी घालणारी वाक्ये,
आरंभीचे स्थितिवर्णन सांगणारी
वाक्ये, तसेच विशिष्ट सूचना
राबवल्यानंतर~$M$ चे स्थितिवर्णन
काय होते ते सांगणारी वाक्ये
असतील.
\end{explain}

\begin{defn}
  \label{tur:und:rep:defn:tm-descr}
$M = \tuple{Q, \Sigma, q_0, \delta}$
हे ट्यूरिंग यंत्र दिले असता
भाषा~$\Lang L_M$ मध्ये पुढील
गोष्टी असतात:
\begin{enumerate}
\item प्रत्येक अवस्था~$q \in Q$ साठी
  द्विस्थानी विधेय $\Obj Q_q(x, y)$.
  अंतर्ज्ञानाच्या दृष्टीने,
  $\Obj Q_q(\num{m}, \num{n})$
  याचा अर्थ “$n$ पायऱ्यांनंतर
  $M$ हे~$q$ अवस्थेत असून
  $m$-व्या घराकडे पाहत आहे.”
\item प्रत्येक चिन्ह~$\sigma\in \Sigma$
  साठी द्विस्थानी विधेय
  $\Obj S_\sigma(x, y)$.
  अंतर्ज्ञानाच्या दृष्टीने,
  $\Obj S_\sigma(\num{m},
  \num{n})$ याचा अर्थ “$n$
  पायऱ्यांनंतर $m$-व्या घरात
  चिन्ह~$\sigma$ आहे.”
\item एक स्थिरांक $\Obj 0$
\item एक एकस्थानी फलन $\prime$
\item एक द्विस्थानी विधेय $<$
\end{enumerate}
\end{defn}

आदान $w = \sigma_{i_1}\dots\sigma_{i_k}$
वर ट्यूरिंग यंत्र~$M$ कसे चालते याचे
वर्णन करणारी वाक्ये
पुढीलप्रमाणे आहेत:
\begin{enumerate}
\item संख्या आणि~$<$ यांचे वर्णन
करणारी स्वयंसिद्धके:
\begin{enumerate}
\item प्रत्येक संख्या तिच्या उत्तरवर्तीपेक्षा
लहान आहे असे सांगणारे वाक्य:
\[\lforall[x][x < x']
\]
\item $<$ संक्रमणशील आहे याची
खात्री करणारे वाक्य:
\[\lforall[x][\lforall[y][\lforall[z][
      ((x < y \land y < z) \lif x < z)]]]
\]
\end{enumerate}
\item आदानाच्या प्रारंभिक स्थितिवर्णनाचे
वर्णन करणारी स्वयंसिद्धके:
\begin{enumerate}
\item $0$~पायऱ्यांनंतर---म्हणजे
  यंत्र सुरू होण्यापूर्वी---$M$
  हे आरंभीच्या अवस्था~$q_0$
  मध्ये असून घर~$1$ कडे पाहते:
\[\Obj Q_{q_0}(\num{1}, \num{0})
\]
\item पहिल्या $k+1$ घरांत
  $\TMendtape$, $\sigma_{i_1}$,
  \dots, $\sigma_{i_k}$ ही चिन्हे
  असतात:
\[\Obj S_\TMendtape(\num{0}, \num{0}) \land
\Obj S_{\sigma_{i_1}}(\num{1}, \num{0}) \land
\dots \land
\Obj S_{\sigma_{i_k}}(\num{k}, \num{0})
\]
\item बाकी फीत रिकामी असते:
\[\lforall[x][(\num{k} < x \lif \Obj S_\TMblank(x, \num{0}))]
\]
\end{enumerate}
\item एका स्थितिवर्णनातून पुढच्या
स्थितिवर्णनाकडे जाणाऱ्या संक्रमणांचे
वर्णन करणारी स्वयंसिद्धके:

पुढे $\varphi(x, y)$ हे पुढील
रूपातील सर्व वाक्यांच्या
संयोगाने मिळते असे मानू:
\[\lforall[z][
  (((z < x \lor x < z) \land \Obj S_\sigma(z, y))
  \lif \Obj S_\sigma(z, y'))]
\]
येथे दोन चल मुक्त असल्याने हा स्वतः स्वतंत्र वाक्य
नसून सूत्रांचा रूपबंध आहे; पुढील संक्रमण-वाक्यांत
हे चल सार्वत्रिक संख्यापकांखाली येतात.

येथे $\sigma \in \Sigma$.
$\varphi(\num{m},\num{n})$ वापरून
“घर~$m$ सोडल्यास, $n+1$
पायऱ्यांनंतरची फीत $n$
पायऱ्यांनंतर होती तशीच आहे”
हे व्यक्त करतो.
\begin{enumerate}
\item \label{tur:und:rep:rep-right} प्रत्येक
  $\delta(q_i, \sigma) =
  \tuple{q_j, \sigma', \TMright}$
  सूचनेसाठी पुढील वाक्य:
\begin{align*}
& \lforall[x][\lforall[y][(
   (\Obj Q_{q_i}(x, y) \land \Obj S_{\sigma}(x, y)) \lif {}]] \\
&\qquad   (\Obj Q_{q_j}(x', y') \land \Obj S_{\sigma'}(x, y') \land
\varphi(x, y)))
\end{align*}
याचा अर्थ असा: $y$ पायऱ्यांनंतर
यंत्र~$q_i$ अवस्थेत असून
चिन्ह~$\sigma$ असलेल्या
घर~$x$ कडे पाहत असेल,
तर $y+1$ पायऱ्यांनंतर
ते घर~$x+1$ कडे पाहील,
~$q_j$ अवस्थेत असेल,
घर~$x$ मध्ये आता~$\sigma'$
असेल आणि घर~$x$ वगळता
इतर प्रत्येक घरात~$y$
पायऱ्यांनंतर होते तेच
चिन्ह असेल.

\item \label{tur:und:rep:rep-left} प्रत्येक
  $\delta(q_i, \sigma) =
  \tuple{q_j, \sigma', \TMleft}$
  सूचनेसाठी पुढील वाक्य:
\begin{align*}
& \lforall[x][\lforall[y][
    ((\Obj Q_{q_i}(x', y) \land \Obj S_{\sigma}(x', y)) \lif {}]]\\
& \qquad   (\Obj Q_{q_j}(x, y') \land \Obj S_{\sigma'}(x', y') \land
\varphi(x', y))) \land {}\\
& \lforall[y][((\Obj Q_{q_i}(\num{0}, y) \land \Obj S_{\sigma}(\num{0},
    y)) \lif {}]\\
& \qquad (\Obj Q_{q_j}(\num{0}, y') \land \Obj S_{\sigma'}(\num{0},
  y') \land \varphi(\num{0}, y)))
\end{align*}
ही रचना नीट समजून घ्या:
आता आपण “घर~$x$ कडे
पाहत असेल तर~\dots”
याऐवजी “घर $x+1$
कडे पाहत असेल तर~\dots”
असे म्हणतो. डावीकडे
जाण्याचा अर्थ पुढच्या
पायरीत यंत्र घर~$x$
कडे पाहते. पण लिहिले
जाते ते घर~$x+1$.
वजाबाकी किंवा पूर्ववर्ती
फलन आपल्याकडे नसल्यामुळे
ही रचना वापरतो.
पहिल्या अटीतील स्थिरता-नियम
लिहिले जाणारे घरच
वगळतो; डावीकडे गेलेले
नवे शीर्षस्थान नाही.


$x+1$ या रूपातील संख्या
$1$, $2$, \dots,
असतात हे लक्षात घ्या.
म्हणून घर~$0$ कडे पाहताना
डावीकडे जाण्याची सूचना
असल्याची बाब यात येत
नाही (आणि अर्थातच ते
डावीकडे जाऊ शकत नाही---त्याच
घरावर राहते). दुसऱ्या
संयुक्त अटीत ही विशेष
बाब येते: $y$ पायऱ्यांनंतर
यंत्र घर~$0$ कडे पाहत,
अवस्था $q_i$ मध्ये असेल
आणि घर~$0$ मध्ये
चिन्ह~$\sigma$ असेल, तर
$y+1$ पायऱ्यांनंतर ते
घर~$0$ कडेच पाहत असेल,
नव्या अवस्था~$q_j$ मध्ये
असेल, घर~$0$ वरील
चिन्ह $\sigma'$ असेल
आणि घर~$0$ वगळता
इतर घरांत $y$~पायऱ्यांनंतर
होती तीच चिन्हे असतील.
\item \label{tur:und:rep:rep-stay} प्रत्येक
  $\delta(q_i, \sigma) =
  \tuple{q_j, \sigma', \TMstay}$
  सूचनेसाठी पुढील वाक्य:
\begin{align*}
& \lforall[x][\lforall[y][(
   (\Obj Q_{q_i}(x, y) \land \Obj S_{\sigma}(x, y)) \lif {}]] \\
&\qquad   (\Obj Q_{q_j}(x, y') \land \Obj S_{\sigma'}(x, y') \land
\varphi(x, y)))
\end{align*}
\end{enumerate}
\end{enumerate}
ट्यूरिंग यंत्र~$M$ आणि आदान~$w$
यांच्यासाठी वरील सर्व वाक्यांच्या संयोगाला $\mathsf{T}(M, w)$ म्हणू.

~$M$ अखेरीस थांबते हे व्यक्त
करण्यासाठी “काही पायऱ्यांनंतर
संक्रमण फलन अपरिभाषित होईल”
असे सांगणारे वाक्य हवे. ज्या सर्व
$\tuple{q, \sigma}$ जोड्यांसाठी
~$\delta(q, \sigma)$ अपरिभाषित
आहे, त्यांचा संच~$X$
असू द्या. मग~$\mathsf{E}(M, w)$
हे पुढील वाक्य
असू द्या:
\[\lexists[x][\lexists[y][(\bigvee_{\tuple{q, \sigma} \in
      X}(\Obj Q_q(x, y) \land \Obj S_\sigma(x, y)))]]
\]

खास थांबण्याची अवस्था~$h$
असलेले ट्यूरिंग यंत्र
वापरले तर हे आणखी सोपे
होते: अशा वेळी पुढील
वाक्य~$\mathsf{E}(M, w)$
यंत्र अखेरीस थांबते हे
व्यक्त करते:
\[\lexists[x][\lexists[y][\Obj Q_h(x, y)]]
\]

\begin{prop}
\label{tur:und:rep:prop:mlessk}
$m < k$ असेल, तर
$\mathsf{T}(M, w) \Entails \num{m} < \num{k}$.
\end{prop}

\begin{proof}
सराव.
\end{proof}

\begin{prob}
\ref{tur:und:rep:prop:mlessk}
सिद्ध करा. (सूचना: $k-m$
वर गणितीय विगमन वापरा.)
\end{prob}











\section{निरूपणाची पडताळणी}\label{tur:und:ver:sec}

\begin{explain}
आपले निरूपण योग्य आहे हे तपासण्यासाठी दोन
गोष्टी सिद्ध कराव्या लागतील. प्रथम,
आदान~$w$ वर $M$ थांबत असेल, तर
$\mathsf{T}(M, w) \lif \mathsf{E}(M, w)$ वैध आहे
हे दाखवायचे. त्यानंतर उलट दिशाही
दाखवायची: $\mathsf{T}(M, w) \lif \mathsf{E}(M, w)$
वैध असेल, तर आदान~$w$ वर
चालवलेले $M$ खरोखर अखेरीस
थांबते.

या दोन सिद्धतांच्या युक्त्या खूप
वेगळ्या आहेत. पहिल्या फलितासाठी
प्रथम-क्रम तर्कशास्त्रातील
वाक्य, म्हणजे
$\mathsf{T}(M, w) \lif \mathsf{E}(M, w)$,
वैध आहे हे दाखवायचे. याचा
सर्वांत सोपा मार्ग म्हणजे
निष्पत्ती देणे.
मात्र आपला पुरावा सर्व $M$
आणि $w$ साठी लागू हवा;
म्हणून ज्या एका विशिष्ट
वाक्याची
निष्पत्ती
द्यायची आहे असे नाही,
तर अशी अनंत वाक्ये आहेत.
त्यामुळे $M$ आणि~$w$
कसेही असले आणि आदान~$w$
वर $M$ थांबायला कितीही
पायऱ्या लागत असल्या,
तरी $\mathsf{T}(M, w) \lif \mathsf{E}(M, w)$
याची निष्पत्ती
मिळेल, हे गणितीय विगमनाने
सिद्ध करणे हा योग्य मार्ग आहे.

अर्थातच, या विगमनाचा आधार
$M$ थांबणाऱ्या स्थितिवर्णनापर्यंत
पोहोचेपर्यंत केलेल्या पायऱ्यांची
संख्या असेल. विगमनात्मक पुराव्यात
आदान~$w$ वरील $M$ च्या
चालक्रमातील प्रत्येक पायरी~$n$
साठी $\mathsf{T}(M, w) \Entails \chi(M, w, n)$
हे सिद्ध करू. येथे $\chi(M, w, n)$
हे आदान~$w$ वर चालवलेल्या~$M$
चे $n$ पायऱ्यांनंतरचे स्थितिवर्णन
अचूक सांगते. आता $M$ आदान~$w$
वर, समजा, $n$ पायऱ्यांनंतर
थांबत असेल, तर $\chi(M, w, n)$
हे थांबणारे स्थितिवर्णन सांगेल.
असे स्थितिवर्णन सांगणारे
$\chi(M, w, n)$ असेल तेव्हा
$\chi(M, w, n) \Entails \mathsf{E}(M, w)$
हेही दाखवू. म्हणून $M$
आदान~$w$ वर थांबत असेल,
तर एखाद्या~$n$ साठी
$n$~पायऱ्यांनंतर $M$
थांबणाऱ्या स्थितिवर्णनात असेल.
त्या वेळी $\mathsf{T}(M, w)
\Entails \chi(M, w, n)$,
आणि $\chi(M, w, n)$
थांबणारे स्थितिवर्णन सांगत
असल्यामुळे $\chi(M, w, n) \Entails \mathsf{E}(M, w)$.
त्यामुळे $\mathsf{T}(M, w) \Entails \mathsf{E}(M, w)$,
म्हणजेच $\Entails \mathsf{T}(M, w)
\lif \mathsf{E}(M, w)$ मिळते.

उलट दिशेची युक्ती फार वेगळी
आहे. येथे $\Entails \mathsf{T}(M, w) \lif \mathsf{E}(M, w)$
असे गृहीत धरून $M$ आदान~$w$
वर थांबते हे सिद्ध करायचे.
गृहीतकावरून $\mathsf{T}(M, w) \Entails \mathsf{E}(M,
w)$ मिळते; म्हणजे $\mathsf{T}(M,
w)$ सत्य असलेल्या प्रत्येक
रचनेत
$\mathsf{E}(M, w)$ सत्य आहे.
म्हणून $\mathsf{T}(M, w)$ सत्य
असलेली एक रचना
~$\Struct{M}$ तयार करू:
तिचे परिभाषाक्षेत्र $\Nat$
असेल, आणि सर्व $\Obj Q_q$
व $\Obj S_\sigma$ यांचे
अर्थनिर्धारण आदान~$w$
वरील~$M$ च्या चालक्रमातील
स्थितिवर्णनांनी ठरेल.
उदाहरणार्थ,
$\Sat{M}{\Obj Q_q(\num{m}, \num{n})}$
तेव्हा आणि केवळ तेव्हाच, जेव्हा $M$ हे
आदान~$w$ वर $n$ पायऱ्या
चालवल्यानंतर अवस्था~$q$
मध्ये असून घर~$m$ कडे
पाहत असेल. आता गृहीतकानुसार
$\mathsf{T}(M, w) \Entails \mathsf{E}(M, w)$
आणि रचनेवरून
$\Sat{M}{\mathsf{T}(M,
  w)}$, म्हणून
$\Sat{M}{\mathsf{E}(M, w)}$.
पण $\Sat{M}{\mathsf{E}(M, w)}$
तेव्हा आणि केवळ तेव्हाच, जेव्हा एखाद्या
$n \in \Domain{M} = \Nat$
साठी आदान~$w$ वर
चालवलेले $M$ हे
~$n$ पायऱ्यांनंतर
थांबणाऱ्या स्थितिवर्णनात
असेल.

\end{explain}

\begin{defn}
$\chi(M, w, n)$ हे पुढील
वाक्य असू द्या:
\[\Obj Q_q(\num{m}, \num{n}) \land \Obj S_{\sigma_0}(\num{0}, \num{n})
\land \dots \land \Obj S_{\sigma_k}(\num{k}, \num{n}) \land
\lforall[x][(\num{k} < x \lif \Obj S_\TMblank(x, \num{n}))]
\]
येथे वेळ~$n$ ला~$q$
ही $M$ ची अवस्था आहे,
वेळ~$n$ ला $M$ घर~$m$
कडे पाहत आहे, वेळ~$n$
ला $0 \le i \le k$
साठी घर~$i$ मध्ये
चिन्ह~$\sigma_i$ आहे,
आणि $k$ हे वेळ~$0$
ला दिलेल्या संपूर्ण आदान-शब्दाच्या शेवटच्या
घराचा क्रमांक किंवा $n$ पायऱ्यांपर्यंत
शीर्षाने भेट दिलेल्या सर्वांत उजव्या
घराचा क्रमांक, यांपैकी अधिक मोठा क्रमांक आहे.
रिकाम्या आदानासाठी शब्दाच्या शेवटचा क्रमांक
शून्य मानतो; आरंभी शीर्ष पहिल्या घरावर असल्याने
ते घरही या मर्यादेत येते.

\end{defn}

\begin{lem}\label{tur:und:ver:lem:halt-config-implies-halt}
आदान~$w$ वर चालवलेले $M$
हे $n$ पायऱ्यांनंतर
थांबणाऱ्या स्थितिवर्णनात
असेल, तर $\chi(M, w, n) \Entails \mathsf{E}(M, w)$.
\end{lem}

\begin{proof}
आदान~$w$ वर $M$ हे $n$
पायऱ्यांनंतर थांबते असे
समजा. मग अशी अवस्था~$q$,
घर~$m$ आणि चिन्ह~$\sigma$
असतील की
\begin{enumerate}
\item $n$ पायऱ्यांनंतर $M$
  हे~$q$ अवस्थेत असून,
  ज्या घर~$m$ वर~$\sigma$
  आहे त्याकडे पाहत असेल.
\item संक्रमण फलन $\delta(q, \sigma)$
  अपरिभाषित असेल.
\end{enumerate}
$\chi(M, w, n)$ हे या
स्थितिवर्णनाचे वर्णन आहे;
त्यात $\Obj Q_{q}(\num{m}, \num{n})$
आणि $\Obj
S_{\sigma}(\num{m}, \num{n})$
ही संयुक्त विधाने असतील.
त्यांवरून एकत्रितपणे
$\mathsf{E}(M,
w)$ मिळते:
\[\lexists[x][\lexists[y][(\bigvee_{\tuple{q, \sigma} \in
      X}(\Obj Q_q(x, y) \land \Obj S_\sigma(x, y)))]]
\]
कारण $\tuple{q,\sigma} \in X$
असल्याने
$\Obj Q_{q}(\num{m}, \num{n}) \land \Obj S_{\sigma}(\num{m},
\num{n}) \Entails \bigvee_{\tuple{q, \sigma} \in X} (\Obj Q_q(\num{m},
\num{n}) \land \Obj S_{\sigma}(\num{m}, \num{n}))$.

\end{proof}

\begin{explain}
म्हणून $M$ आदान $w$
वर थांबत असेल, तर
एखाद्या~$n$ साठी
$\chi(M,
w, n) \Entails \mathsf{E}(M,w)$.
आता चालक्रम पोहोचतो अशा कोणत्याही वेळ~$n$
साठी $\mathsf{T}(M, w) \Entails \chi(M, w, n)$
हे दाखवू.
\end{explain}

\begin{lem}
\label{tur:und:ver:lem:config}
प्रत्येक $n$ साठी,
$M$ चा चालक्रम $n$ पायऱ्यांनंतरच्या
स्थितिवर्णनापर्यंत पोहोचत असेल, तर
$\mathsf{T}(M, w)
\Entails \chi(M, w, n)$.
\end{lem}


\begin{proof}
विगमनाची आरंभीची पायरी:
$n = 0$ असेल, तर
$\chi(M, w, 0)$ मधील अवस्था व फितीची विधाने
$\mathsf{T}(M, w)$ च्या आरंभीच्या अटींवरून मिळतात.
आदान रिकामे नसल्यास ही त्यातीलच संयुक्त विधाने आहेत.
आदान रिकामे असल्यास शीर्ष पहिल्या घरावर आहे;
तेथील रिकामे चिन्ह शून्यापेक्षा त्याचा क्रमांक मोठा
असल्याने मिळते. पहिल्या घरापलीकडील रिकामेपण
क्रमाच्या संक्रमणशीलतेने शून्यापलीकडील
रिकामेपणावरून मिळते.

विगमनाचे गृहीतक:
$n$-व्या पायरीपूर्वी $M$
थांबलेले नसेल, तर
$\mathsf{T}(M,w) \Entails \chi(M, w, n)$.
आता $\chi(M, w, n)$
थांबणारे स्थितिवर्णन सांगत
नसेल, तर $\mathsf{T}(M, w) \Entails
\chi(M, w, n+1)$ हे
दाखवायचे आहे.

$n \ge 0$ असे समजा
आणि आदान $w$ वर
सुरू केलेले $M$ हे
$n$ पायऱ्यांनंतर
अवस्था~$q$ मध्ये असून
घर~$m$ कडे पाहत
असेल. $n$ पायऱ्यांनंतर
$M$ थांबत नसल्यामुळे
$M$ च्या कार्यक्रमात
पुढील तीनपैकी एका
रूपाची सूचना असली पाहिजे:

\begin{enumerate}
\item \label{tur:und:ver:right} $\delta(q, \sigma) = \tuple{q', \sigma', \TMright}$

\item \label{tur:und:ver:left} $\delta(q, \sigma) = \tuple{q', \sigma', \TMleft}$

\item \label{tur:und:ver:stay} $\delta(q, \sigma) = \tuple{q', \sigma', \TMstay}$
\end{enumerate}

या तीनही बाबींचा आता
एकेक करून विचार करू.

\begin{enumerate}
\item \ref{tur:und:ver:right} या रूपाची सूचना
  आहे असे समजा. अशा वेळी
  \ref{tur:und:rep:defn:tm-descr}\ref{tur:und:rep:rep-right}
  यांनुसार
\begin{align*}
& \lforall[x][\lforall[y][((\Obj Q_{q}(x,
  y) \land \Obj S_\sigma(x, y)) \lif {}]]\\
& \qquad (\Obj Q_{q'}(x',
  y') \land \Obj S_{\sigma'}(x, y') \land \varphi(x, y)))
\intertext{हे $\mathsf{T}(M,w)$ चे एक संयुक्त विधान आहे.
  त्यावरून पुढील वाक्य
  मिळते (सार्वत्रिक विधानात~$x$ ऐवजी
  $\num{m}$ आणि~$y$ ऐवजी
  $\num{n}$ बसवून):}
& (\Obj Q_{q}(\num{m}, \num{n}) \land \Obj S_{\sigma}(\num{m},
\num{n})) \lif {}\\
& \qquad (\Obj Q_{q'}(\num{m}', \num{n}') \land
\Obj S_{\sigma'}(\num{m}, \num{n}') \land \varphi(\num{m}, \num{n})).
\intertext{विगमनाच्या गृहीतकानुसार
  $\mathsf{T}(M, w) \Entails \chi(M, w, n)$,
  म्हणजे}
& \Obj Q_q(\num{m}, \num{n}) \land \Obj S_{\sigma_0}(\num{0}, \num{n})
\land \dots \land \Obj S_{\sigma_k}(\num{k}, \num{n}) \land \\
& \qquad\lforall[x][(\num{k} < x \lif \Obj S_\TMblank(x, \num{n}))]\\
\intertext{$n$ पायऱ्यांनंतर फितीवरील घर~$m$ मध्ये
  चिन्ह~$\sigma$ असल्याने, त्याच्याशी संबंधित
  संयुक्त विधान~$\Obj S_\sigma(\num{m}, \num{n})$
  आहे; म्हणून यावरून पुढील मिळते:}
& \Obj Q_{q}(\num{m}, \num{n}) \land \Obj S_{\sigma}(\num{m},
   \num{n})
\intertext{आता आपल्याला पुढील मिळते:}
& \Obj Q_{q'}(\num{m}', \num{n}') \land \Obj S_{\sigma'}(\num{m},
  \num{n}') \land {}\\
& \qquad\Obj S_{\sigma_0^{+}}(\num{0}, \num{n}') \land \dots \land
  \Obj S_{\sigma_k^{+}}(\num{k}, \num{n}') \land {}\\
& \qquad \lforall[x][(\num{k} < x \lif \Obj S_\TMblank(x, \num{n}'))]
\end{align*}
येथे वरची अधिक खूण पुढील स्थितीतील चिन्ह
दर्शवते: लिहिलेल्या घरात नव्या सूचनेतील चिन्ह,
आणि इतर घरांत आधीची चिन्हे आहेत.
याचे कारण असे: पहिली ओळ आधीच्या
सोपाधिक विधानाच्या परिणामी भागावरून
निगमन नियमाने थेट मिळते.
मधल्या ओळीत लिहिलेल्या घरातील नवे चिन्ह
पहिल्या ओळीवरूनच मिळते; जुन्या चिन्हाचे
$S_{\sigma_m}(\num{m},\num{n}')$
हे विधान राखलेले नाही. इतर घरांतील प्रत्येक
विधान हे~$\chi(M, w, n)$ मधील संबंधित
संयुक्त विधान आणि $\varphi(\num{m},
\num{n})$ यांवरून मिळते.

$m < k$ असेल, तर
$\mathsf{T}(M,w) \Proves \num{m} < \num {k}$
(\ref{tur:und:rep:prop:mlessk}),
आणि~$<$ च्या संक्रमणशीलतेवरून
$\lforall[x][(\num{k} < x \lif \num{m} < x)]$
मिळते. $m = k$ असेल,
तर $\lforall[x][(\num{k} < x \lif \num{m} < x)]$
केवळ तर्कनियमांवरून मिळते.
म्हणून शेवटची ओळ
$\chi(M, w,
n)$ मधील संबंधित संयुक्त विधान,
$\lforall[x][(\num{k} < x \lif \num{m} < x)]$,
आणि $\varphi(\num{m},
\num{n})$ यांवरून मिळते.
$m<k$ असेल, तर हेच
$\chi(M, w, n+1)$ आहे.

आता $m=k$ असे समजा.
या वेळी $n+1$ पायऱ्यांनंतर
शीर्षाने घर~$k+1$ लाही
भेट दिलेली असेल; तेच
आता भेट दिलेले सर्वांत
उजवे घर आहे. म्हणून
$\chi(M, w, n+1)$ मध्ये
$\Obj
S_\TMblank(\num{k}',\num{n}')$
हे नवे संयुक्त विधान असेल,
आणि शेवटचे संयुक्त विधान
$\lforall[x][(\num{k}' < x \lif \Obj S_\TMblank(x, \num{n}'))]$
असेल. ही दोन्ही वाक्ये
देखील निष्पन्न होतात
हे तपासावे लागेल.

आपल्याकडे आधीच
$\lforall[x][(\num{k} < x \lif \Obj S_\TMblank(x,
  \num{n}'))]$
आहे. विशेषतः त्यावरून
$\num{k} < \num{k}' \lif
\Obj S_\TMblank(\num{k}', \num{n}')$
मिळते. स्वयंसिद्धक
$\lforall[x][x <
  x']$ यावरून
$\num{k} < \num{k}'$
मिळते. सोपाधिक
निगमनाने
$\Obj
S_\TMblank(\num{k}',\num{n}')$
मिळते.

शिवाय,
$\mathsf{T}(M,w) \Proves \num{k} < \num{k}'$
असल्यामुळे~$<$ च्या
संक्रमणशीलतेच्या
स्वयंसिद्धकावरून
$\lforall[x][(\num{k}' < x \lif \Obj
  S_\TMblank(x, \num{n}'))]$
मिळते. (याची पडताळणी
सरावासाठी सोडतो.)

\item \ref{tur:und:ver:left} या रूपाची
सूचना आहे असे समजा.
मग \ref{tur:und:rep:defn:tm-descr}\ref{tur:und:rep:rep-left}
यांनुसार
\begin{align*}
& \lforall[x][\lforall[y][((\Obj Q_{q}(x', y) \land \Obj
    S_{\sigma}(x', y)) \lif {}]]\\
& \qquad (\Obj Q_{q'}(x, y') \land \Obj
  S_{\sigma'}(x', y') \land \varphi(x', y))) \land {}\\
& \lforall[y][((\Obj Q_{q_i}(\num{0}, y) \land \Obj S_{\sigma}(\num{0},
    y)) \lif {}]\\
& \qquad (\Obj Q_{q_j}(\num{0}, y') \land \Obj S_{\sigma'}(\num{0}, y')
  \land \varphi(\num{0}, y)))
\intertext{हे $\mathsf{T}(M,w)$ चे एक
  संयुक्त विधान आहे. $m>0$
  असेल, तर $l = m - 1$
  असे घ्या (म्हणजे
  $m = l+1$). वरील
  वाक्याच्या
  पहिल्या संयुक्त विधानावरून
  पुढील मिळते:}
& (\Obj Q_{q}(\num{l}', \num{n}) \land \Obj S_{\sigma}(\num{l}', \num{n}))
\lif {} \\
& \qquad (\Obj Q_{q'}(\num{l}, \num{n}') \land \Obj S_{\sigma'}(\num{l}',
\num{n}') \land \varphi(\num{l}', \num{n}))
\intertext{अन्यथा $l = m = 0$
  असे घ्या आणि दुसऱ्या
  संयुक्त विधानावरून निष्पन्न
  होणारे पुढील वाक्य पाहा:}
& ((\Obj Q_{q_i}(\num{0}, \num{n}) \land \Obj S_{\sigma}(\num{0}, \num{n})) \lif {}\\
& \qquad   (\Obj Q_{q_j}(\num{0}, \num{n}') \land \Obj S_{\sigma'}(\num{0}, \num{n}') \land
\varphi(\num{0}, \num{n})))
\intertext{शून्यावरील शाखेतील पहिली अवस्था
  आधीची, तर दुसरी पुढची अवस्था आहे.
  दोन्हीपैकी कोणत्याही
  वाक्यावरून पुढील मिळते:}
&  \Obj Q_{q'}(\num{l}, \num{n}') \land \Obj S_{\sigma'}(\num{m},
  \num{n}') \land {}\\
&\qquad  \Obj S_{\sigma_0^{+}}(\num{0}, \num{n}')\land \dots \land
  \Obj S_{\sigma_k^{+}}(\num{k}, \num{n}')  \land {} \\
& \qquad  \lforall[x][(\num{k} < x
  \lif \Obj S_\TMblank(x, \num{n}'))]
\end{align*}
पहिल्या शाखेत स्थिरता-नियम लिहिले गेलेले
मूळ घर वगळतो; डावीकडे गेलेले नवे शीर्षस्थान नाही.
येथेही वरची अधिक खूण पुढील स्थितीतील
चिन्ह दर्शवते; लिहिलेल्या घरात नवे चिन्ह
आणि इतर घरांत आधीची चिन्हे आहेत.


आधीप्रमाणेच. (पहिल्या बाबतीत
$\num{l}' \ident \num{l+1}
\ident \num{m}$ आणि
दुसऱ्या बाबतीत
$\num{l} \ident \num{0}$
हे लक्षात घ्या.) पण
हेच $\chi(M, w, n+1)$
आहे.

\item \ref{tur:und:ver:stay} ही बाब
सरावासाठी सोडतो.
\end{enumerate}
म्हणून चालक्रम पोहोचतो अशा कोणत्याही~$n$
साठी $\mathsf{T}(M, w) \Entails \chi(M, w, n)$
हे दाखवले.
\end{proof}

\begin{prob}
\ref{tur:und:ver:lem:config}
च्या पुराव्यातील
~\ref{tur:und:ver:stay}
ही बाब पूर्ण करा.
\end{prob}

\begin{prob}
$\Obj S_{\sigma_i}(\num{i}, \num{n})$
आणि $\varphi(m, n)$ यांवरून
$\Obj S_{\sigma_i}(\num{i}, \num{n}')$
ची निष्पत्ती द्या
($i \neq
m$, म्हणजे $i < m$
किंवा $m < i$,
असे गृहीत धरून).
\end{prob}

\begin{prob}
$\lforall[x][(\num{k} < x \lif \Obj
  S_\TMblank(x, \num{n}'))]$,
$\lforall[x][x < x']$,
आणि
$\lforall[x][\lforall[y][\lforall[z][ ((x < y \land y < z) \lif x <
      z)]]]$
यांवरून
$\lforall[x][(\num{k}' < x \lif \Obj
  S_\TMblank(x, \num{n}'))]$
ची निष्पत्ती द्या.
\end{prob}


\begin{lem}
\label{tur:und:ver:lem:valid-if-halt}
$M$ आदान~$w$ वर
थांबत असेल, तर
$\mathsf{T}(M, w) \lif
\mathsf{E}(M, w)$ वैध आहे.
\end{lem}

\begin{proof}
\ref{tur:und:ver:lem:config} नुसार,
चालक्रम पोहोचतो अशा कोणत्याही वेळ~$n$
ला~$M$ चे वेळ~$n$ वरील स्थितिवर्णन
सांगणारे~$\chi(M, w, n)$
हे~$\mathsf{T}(M, w)$ वरून
निष्पन्न होते. समजा,
$M$ हे $k$ पायऱ्यांनंतर
थांबते. त्या वेळी
ते कोणत्यातरी~$m \in \Nat$
साठी घर~$m$ कडे
पाहत असेल. मग
$\chi(M, w, k)$ हे~$M$
चे थांबणारे स्थितिवर्णन
सांगते. म्हणजे त्यात
$\Obj Q_q(\num{m}, \num{k})$
आणि $\Obj
S_\sigma(\num{m}, \num{k})$
ही दोन्ही संयुक्त विधाने
असतील व $\delta(q,\sigma)$
अपरिभाषित असेल. म्हणून
\ref{tur:und:ver:lem:halt-config-implies-halt}
नुसार
$\chi(M, w, k) \Entails \mathsf{E}(M,
w)$. पण
$\mathsf{T}(M, w) \Entails \chi(M, w, k)$
असल्यामुळे
$\mathsf{T}(M, w)
\Entails \mathsf{E}(M, w)$;
म्हणून $\mathsf{T}(M, w) \lif \mathsf{E}(M, w)$
वैध आहे.
\end{proof}


\begin{explain}
आपला दावा पूर्णपणे
पडताळण्यासाठी उलटी
दिशाही स्थापित करावी
लागेल: $\mathsf{T}(M, w) \lif \mathsf{E}(M, w)$
वैध असेल, तर
आदान~$w$ वर सुरू
केलेले $M$ खरोखर
थांबते.
\end{explain}

\begin{lem}
\label{tur:und:ver:lem:halt-if-valid}
$\Entails \mathsf{T}(M, w) \lif \mathsf{E}(M, w)$
असेल, तर $M$
आदान~$w$ वर थांबते.
\end{lem}

\begin{proof}
परिभाषाक्षेत्र~$\Nat$ असलेली
$\Lang L_M$-रचना
~$\Struct M$ घ्या. ती
$\Obj 0$ चा अर्थ~$0$,
$\prime$ चा अर्थ उत्तरवर्ती
फलन आणि $<$ चा अर्थ
लहान-असणे संबंध असा
लावते. $\Obj Q_q$
आणि~$\Obj S_\sigma$
ही विधेये पुढीलप्रमाणे
अर्थनिर्धारित करतात:
\begin{align*}
  \Assign{\Obj Q_q}{M} & =
\Setabs{\tuple{m, n}}{\begin{array}{ll}\text{आदान~$w$ वर सुरू होऊन $n$ पायऱ्यांनंतर,}\\ \text{$M$ हे अवस्था $q$
  मध्ये असून घर~$m$ कडे पाहते}\end{array}} \\
\Assign{\Obj S_\sigma}{M} & = \Setabs{\tuple{m, n}}{\begin{array}{ll}
\text{आदान~$w$ वर सुरू होऊन $n$ पायऱ्यांनंतर,}\\ \text{$M$ च्या घर~$m$ मध्ये
  चिन्ह~$\sigma$ आहे}\end{array}}
\end{align*}
म्हणजे, आदान~$w$ वर
सुरू केलेले $M$ प्रत्यक्षात
प्रत्येक पायरीला काय करते
हे वर्णन करणारी रचना~$\Struct{M}$
आपण बांधतो. स्पष्टपणे,
$\Sat{M}{\mathsf{T}(M, w)}$.
$\Entails \mathsf{T}(M, w) \lif \mathsf{E}(M, w)$
असेल, तर
$\Sat{M}{\mathsf{E}(M, w)}$
देखील; म्हणजे
\[\Sat{M}{\lexists[x][\lexists[y][(\bigvee_{\tuple{q, \sigma} \in
      X}(\Obj Q_q(x, y) \land \Obj S_\sigma(x, y)))]]}.
\]
$\Domain{M} = \Nat$
असल्यामुळे अशी
$m$, $n \in \Nat$
असतील की काही~$q$
आणि~$\sigma$ साठी
$\delta(q, \sigma)$
अपरिभाषित असून
$\Sat{M}{\Obj Q_q(\num{m}, \num{n}) \land \Obj S_\sigma(\num{m},
\num{n})}$
असेल. ~$\Struct M$
च्या व्याख्येनुसार,
याचा अर्थ आदान~$w$
वर सुरू केलेले $M$
हे~$n$ पायऱ्यांनंतर
अवस्था~$q$ मध्ये असून
चिन्ह~$\sigma$ वाचत आहे,
आणि संक्रमण फलन
अपरिभाषित आहे;
म्हणजे~$M$ थांबले आहे.
\end{proof}











\section{निर्णय समस्या सोडवता येत नाही}\label{tur:und:uns:sec}

\begin{thm}
\label{tur:und:uns:thm:decision-prob}
निर्णय समस्या सोडवता येत नाही:
फितीवर प्रथम-क्रम तर्कशास्त्रातील
वाक्य~$\psi$ आदान म्हणून
असताना सुरू केले की $D$ अखेरीस थांबून,
$\psi$ वैध असेल तेव्हा आणि केवळ तेव्हाच~$1$ आणि
अन्यथा~$0$ फलित देणारे
कोणतेही ट्यूरिंग यंत्र~$D$ नाही.
\end{thm}

\begin{proof}
निर्णय समस्या सोडवता येते,
म्हणजे असे ट्यूरिंग यंत्र~$D$
आहे, असे समजा. मग पुढीलप्रमाणे
थांबण्याची समस्या सोडवता येईल.
ट्यूरिंग यंत्र~$M_e$ चा संकेतांक~$e$
आणि आदान~$w$ दिले असता
संबंधित वाक्य
$\mathsf{T}(M_e, w) \lif \mathsf{E}(M_e, w)$
संगणित करून अंतखूणेच्या लगेच उजवीकडील
पहिल्या आदानघराकडे पाहत थांबणारे
ट्यूरिंग यंत्र~$E$ बांधू.

मग आदान~$e$ आणि~$w$
दिले असता $E \concat D$
प्रथम $\mathsf{T}(M_e, w) \lif
\mathsf{E}(M_e, w)$ संगणित करेल,
आणि नंतर त्या आदानावर
निर्णय-समस्येचे यंत्र~$D$
चालवेल. $\mathsf{T}(M_e, w) \lif \mathsf{E}(M_e, w)$
वैध असेल तेव्हा आणि केवळ तेव्हाच $D$
$1$ देऊन थांबेल, आणि
अन्यथा~$0$ देईल.
\ref{tur:und:ver:lem:halt-if-valid}
आणि \ref{tur:und:ver:lem:valid-if-halt}
यांनुसार $\mathsf{T}(M_e, w) \lif \mathsf{E}(M_e, w)$
वैध असते तेव्हा आणि केवळ तेव्हाच~$M_e$
आदान~$w$ वर थांबते.
म्हणून आदान~$e$ आणि~$w$
दिले असता $E\concat D$
हे~$M_e$ आदान~$w$
वर थांबत असेल तेव्हा आणि केवळ तेव्हाच
$1$ देऊन थांबेल, आणि
अन्यथा~$0$ देऊन थांबेल.
दुसऱ्या शब्दांत, $E \concat D$
थांबण्याची समस्या सोडवेल.
पण \ref{tur:und:hal:thm:halting-problem}
नुसार असे कोणतेही
ट्यूरिंग यंत्र असू शकत नाही.
\end{proof}

\begin{cor}\label{tur:und:uns:cor:undecidable-sat}
प्रथम-क्रम तर्कशास्त्रातील
कोणतेही वाक्य
पूर्ततायोग्य आहे की नाही,
हे ठरवण्याची समस्या अनिर्णेय आहे.
\end{cor}

\begin{proof}
  पूर्ततायोग्यता ट्यूरिंग यंत्र~$S$
  ठरवू शकते असे समजा. मग
  निर्णय समस्या पुढीलप्रमाणे
  सोडवता येईल: आदान म्हणून
  वाक्य~$\psi$ दिले
  असता~$\psi$ एक घर उजवीकडे
  सरकवा. घर~$1$ कडे परत
  या आणि तेथे~$\lnot$
  हे चिन्ह लिहा.


  आता ट्यूरिंग यंत्र~$S$
  चालवा. फितीवर ते अखेरीस
  $1$ ($\lnot \psi$ पूर्ततायोग्य
  असेल तर) किंवा~$0$
  ($\lnot \psi$ अपूर्ततायोग्य
  असेल तर) फलित देऊन
  थांबेल. फलित मिळाल्यावर शीर्ष पहिल्या
  आदानघरावर परत आणा. घर~$1$ वर
  $\TMstroke$ असेल, तर
  तो पुसा; घर~$1$
  रिकामे असेल, तर तेथे
  $\TMstroke$ लिहा.
  त्यानंतर थांबा.

  हे ट्यूरिंग यंत्र नेहमी
  थांबते. $\lnot
  \psi$ अपूर्ततायोग्य असेल
  तेव्हा आणि केवळ तेव्हाच ते~$1$ देते;
  अन्यथा~$0$ देते.
  कारण $\psi$ वैध असते
  तेव्हा आणि केवळ तेव्हाच~$\lnot \psi$
  अपूर्ततायोग्य असते,
  यंत्र~$\psi$ वैध असेल
  तेव्हा आणि केवळ तेव्हाच~$1$ आणि
  अन्यथा~$0$ देईल.
  म्हणजे ते निर्णय
  समस्या सोडवेल.
\end{proof}

\begin{explain}
म्हणून “$\psi$ हे प्रथम-क्रम
तर्कशास्त्रातील वैध वाक्य आहे का?” या प्रश्नाला
नेहमी बरोबर “हो” किंवा
“नाही” उत्तर देणारे
ट्यूरिंग यंत्र नाही. मात्र
“हो” हे उत्तर बरोबर असेल
तेव्हा ते देणारे ट्यूरिंग
यंत्र \emph{आहे}; उत्तर
“नाही” असेल, तर ते
थांबतच नाही. हे प्रथम-क्रम
तर्कशास्त्राच्या निर्दोषता आणि
संपूर्णता प्रमेयांवरून, तसेच
निष्पत्त्या प्रभावी
रीतीने प्रगणित करता येतात
यावरून मिळते.
\end{explain}

\begin{thm}
  \label{tur:und:uns:thm:valid-ce}
  प्रथम-क्रम वाक्यांची वैधता अर्धनिर्णेय आहे:
  फितीवर प्रथम-क्रम
  तर्कशास्त्रातील वाक्य~$\psi$ आदान म्हणून
  असताना सुरू केले की
  $\psi$ वैध असेल तेव्हा आणि केवळ तेव्हाच
  अखेरीस थांबून~$1$
  देणारे, आणि अन्यथा
  न थांबणारे ट्यूरिंग
  यंत्र~$E$ आहे. हे $E$
  यंत्र वैधतेला होकार देते.
\end{thm}

\begin{proof}
  प्रथम-क्रम तर्कशास्त्रातील
  सर्व शक्य निष्पत्त्या
  एका प्रभावी अल्गोरिदमने
  एकामागून एक तयार करता
  येतात. यंत्र~$E$ तसे
  करते; $\Proves
  \psi$ हे दाखवणारी
  निष्पत्ती
  सापडली की ते~$1$
  देऊन थांबते. निर्दोषता
  प्रमेयामुळे $E$~हे $1$
  देऊन थांबले असेल,
  तर~$\Entails \psi$.
  संपूर्णता प्रमेयामुळे
  $\Entails \psi$ असेल,
  तर~$\Proves \psi$
  हे दाखवणारी निष्पत्ती असते. $E$
  सर्व शक्य निष्पत्त्या
  क्रमाने तयार करत असल्याने
  $\Proves \psi$ दाखवणारी
  अशी निष्पत्ती त्याला
  अखेरीस सापडेल; म्हणून
  ते~$1$ देऊन थांबेल.
\end{proof}











\section{ट्राख्टेनब्रोटचे प्रमेय}\label{tur:und:tra:sec}

\begin{explain}
  \ref{tur:und:rep:sec} मध्ये ट्यूरिंग
  यंत्र~$M$ आणि आदान
  चिन्हमाला~$w$ यांच्यासाठी
  $\mathsf{T}(M,w)$ आणि~$\mathsf{E}(M,w)$
  ही वाक्ये
  परिभाषित केली. त्यानंतर
  \ref{tur:und:ver:lem:valid-if-halt}
  आणि \ref{tur:und:ver:lem:halt-if-valid}
  मध्ये आदान~$w$ वर
  सुरू केलेले~$M$ अखेरीस
  थांबते तेव्हा आणि केवळ तेव्हाच
  $\mathsf{T}(M,w) \lif \mathsf{E}(M,w)$
  वैध असते, हे दाखवले.
  थांबण्याची समस्या अनिर्णेय
  असल्यामुळे प्रथम-क्रम
  तर्कशास्त्रातील वाक्यांची वैधता आणि
  पूर्ततायोग्यताही अनिर्णेय आहेत
  (\ref{tur:und:uns:thm:decision-prob} आणि \ref{tur:und:uns:cor:undecidable-sat}).

  पण वाक्यांची वैधता आणि
  पूर्ततायोग्यता कोणत्याही
  रचनांसाठी
  परिभाषित होतात, मग त्या
  सांत असोत वा अनंत.
  एखादे वाक्य
  सांत रचनेत
  पूर्ततायोग्य आहे का
  (किंवा सर्व सांत
  रचनांमध्ये
  वैध आहे का), हे ठरवणे
  सोपे असेल असे वाटू शकते.
  निर्णय समस्या सोडवता
  येत नाही याचा पुरावा
  आपण यासाठीही रूपांतरित
  करू शकतो; तसे नसते
  हे त्यातून दिसेल.

  प्रथम \ref{tur:und:ver:lem:halt-if-valid}
  च्या पुराव्याकडे परत पाहा.
  तेथे आपण~$\mathsf{T}(M,w)$ चे
  प्रतिरूप~$\Struct M$
  तयार केले; आदान~$w$
  वर सुरू केलेले यंत्र~$M$
  नेमके काय करते ते
  त्यात वर्णिले आहे.
  त्या प्रतिरूपाचे परिभाषाक्षेत्र
  $\Nat$, म्हणजे अनंत,
  होते. पण $M$ हे
  आदान~$w$ वर खरोखर
  थांबत असेल, तर त्याच
  पद्धतीने सांत प्रतिरूप
  ~$\Struct M'$ बांधता येते.
  समजा आदान~$w$ वर
  सुरू झालेले $M$ हे
  $k$ पायऱ्यांनंतर थांबते.
  $\Domain{M'}$ हे
  $\{0, \dots, n\}$ घ्या;
  येथे $n$ ही संख्या
  $k$ पेक्षा एकाने मोठी
  संख्या आणि~$w$ च्या
  लांबीपेक्षा एकाने मोठी
  संख्या यांपैकी मोठी
  आहे. तसेच
  \[    \Assign{\prime}{M'}(x) =
    \begin{cases}
      x + 1 &\text{$x < n$ असेल तर}\\
      n &\text{अन्यथा,}
    \end{cases}
  \]
  घ्या, आणि
  $\tuple{x,y} \in \Assign{<}{M'}$
  तेव्हा आणि केवळ तेव्हाच, जेव्हा $x < y$
  किंवा $x = y = n$.
  इतर बाबतीत $\Struct{M'}$
  ची व्याख्या~$\Struct{M}$
  प्रमाणेच आहे. $\Struct{M'}$
  च्या व्याख्येनुसार,
  \ref{tur:und:ver:lem:halt-if-valid}
  च्या पुराव्यातल्याप्रमाणे
  $\Sat{M'}{\mathsf{T}(M,w)}$.
  तसेच $M$ हे आदान~$w$
  वर थांबते असे गृहीत
  धरल्यामुळे
  $\Sat{M'}{\mathsf{E}(M,w)}$.
  म्हणून~$\Struct{M'}$
  हे $\mathsf{T}(M,w) \land \mathsf{E}(M,w)$
  चे सांत प्रतिरूप आहे
  (येथे $\lif$ ऐवजी
  $\land$ वापरले आहे).


  पुराव्याचा अर्धा भाग
  झाला आहे: $M$ हे
  आदान~$w$ वर थांबले,
  तर $\mathsf{T}(M,w) \land \mathsf{E}(M,w)$
  ला सांत प्रतिरूप
  असते, हे दाखवले.
  दुर्दैवाने याचा उलट
  निष्कर्ष खरा नाही:
  काही ट्यूरिंग यंत्रे
  आदान~$w$ वर थांबत
  नाहीत, तरी
  $\mathsf{T}(M,w)
  \land \mathsf{E}(M,w)$ ला
  सांत प्रतिरूप असू शकते.
  उदाहरणार्थ, एकच
  अवस्था $q_0$ आणि
  $\delta(q_0,\TMblank) = \tuple{q_0,\TMblank,\TMstay}$
  ही एकच सूचना
  असलेले यंत्र~$M$
  घ्या. रिकाम्या आदानावर
  $w = \emptyseq$
  सुरू केल्यावर ते
  कधीच थांबत नाही:
  ते अनंत फेरीत
  राहते, पण फीत
  बदलत नाही किंवा
  शीर्ष हलवत नाही.
  सर्व स्थितिवर्णने
  सारखीच असतात
  (तीच अवस्था,
  तेच शीर्षस्थान,
  तीच फीत).
  तरी $\mathsf{T}(M,\emptyseq) \land \mathsf{E}(M,\emptyseq)$
  सत्य करणारी सांत
  रचना
  ~$\Struct{M''}$
  परिभाषित करता येते
  (सराव). पण अशा
  रचना
  वगळण्यासाठी $\mathsf{T}(M,w)$
  योग्य रीतीने बदलता
  येते.

  \begin{prob}
    $q_0$ ही एकच
    अवस्था आणि
    $\delta(q_0,\TMblank) = \tuple{q_0,\TMblank,\TMstay}$
    ही एकच सूचना
    असलेले ट्यूरिंग
    यंत्र~$M$ घ्या.
    $\Domain{M''} = \{0, 1, 2\}$,
    $\Assign{\prime}{M''}(0) =
    \Assign{\prime}{M''}(1) = 1$
    आणि $\Assign{\prime}{M''}(2) = 2$,
    तसेच $\Assign{<}{M''} =
    \{\tuple{0,1}, \tuple{1,1}, \tuple{2,2}\}$
    असे घ्या.
    $\Assign{\Obj Q_{q_0}}{M''}$,
    $\Assign{\Obj
    S_{\TMblank}}{M''}$
    आणि $\Assign{\Obj S_{\TMendtape}}{M''}$
    अशी परिभाषित करा
    की $\mathsf{T}(M,\emptyseq)$
    आणि $\mathsf{E}(M,\emptyseq)$
    दोन्ही सत्य ठरतील;
    तसे का होते ते
    स्पष्ट करा.
    सूचना: $\delta(q_0, \TMendtape)$
    अपरिभाषित आहे हे
    लक्षात घ्या.
    पुढील दोन्ही विधाने
    ~$\Struct{M''}$
    मध्ये सत्य आहेत
    याची खात्री करा:
    \begin{align*}
    & \Obj Q_{q_0}(\num{1}, \num{n}) \land \Obj S_{\TMendtape}(\num{0}, \num{n})
      \land
      \lforall[x][(\num{0} < x \lif \Obj S_\TMblank(x, \num{n}))] \text{\quad सर्व $n \in \Nat$ साठी}\\
    & \lexists[y][(\Obj Q_{q_0}(\num{0}, y) \land \Obj S_{\TMendtape}(\num{0}, y))]
    \end{align*}
  \end{prob}

\end{explain}

ट्यूरिंग यंत्र~$M$ हे
आदान $w = \sigma_{i_1}\dots\sigma_{i_k}$
वर कसे चालते, याचे
वर्णन करणारी पुढील
वाक्ये पाहा:
\begin{enumerate}
\item संख्या आणि~$<$
यांचे वर्णन करणारी
स्वयंसिद्धके (\ref{tur:und:rep:sec}
मधील~$\mathsf{T}(M,w)$ च्या
व्याख्येप्रमाणेच).
\item आरंभीच्या स्थितिवर्णनाचे
वर्णन करणारी स्वयंसिद्धके:
तीही~$\mathsf{T}(M,w)$
च्या व्याख्येप्रमाणेच.
\item एका स्थितिवर्णनातून
पुढच्या स्थितिवर्णनात
होणाऱ्या संक्रमणाचे
वर्णन करणारी स्वयंसिद्धके:

पुढील बाबींसाठी $\varphi(x, y)$
आधीप्रमाणे असू द्या,
आणि
\[  \psi(y) \ident \lforall[x][(x < y \lif \eq/[x][y])].
\]
\begin{enumerate}
\item \label{tur:und:tra:rep-right}
$\delta(q_i, \sigma) =
  \tuple{q_j, \sigma', \TMright}$
रूपाच्या प्रत्येक सूचनेसाठी
पुढील वाक्य:
\begin{align*}
& \lforall[x][\lforall[y][(
   (\Obj Q_{q_i}(x, y) \land \Obj S_{\sigma}(x, y)) \lif {}]] \\
&\qquad   (\Obj Q_{q_j}(x', y') \land \Obj S_{\sigma'}(x, y') \land
\varphi(x, y) \land \psi(y')))
\end{align*}
\item \label{tur:und:tra:rep-left}
$\delta(q_i, \sigma) =
  \tuple{q_j, \sigma', \TMleft}$
रूपाच्या प्रत्येक सूचनेसाठी
पुढील वाक्य:
\begin{align*}
& \lforall[x][\lforall[y][
    ((\Obj Q_{q_i}(x', y) \land \Obj S_{\sigma}(x', y)) \lif {}]]\\
& \qquad   (\Obj Q_{q_j}(x, y') \land \Obj S_{\sigma'}(x', y') \land
\varphi(x', y) \land \psi(y'))) \land {}\\
& \lforall[y][((\Obj Q_{q_i}(\num{0}, y) \land \Obj S_{\sigma}(\num{0},
    y)) \lif {}]\\
& \qquad (\Obj Q_{q_j}(\num{0}, y') \land \Obj S_{\sigma'}(\num{0},
  y') \land \varphi(\num{0}, y) \land \psi(y')))
\end{align*}
\item \label{tur:und:tra:rep-stay}
$\delta(q_i, \sigma) =
  \tuple{q_j, \sigma', \TMstay}$
रूपाच्या प्रत्येक सूचनेसाठी
पुढील वाक्य:
\begin{align*}
& \lforall[x][\lforall[y][(
   (\Obj Q_{q_i}(x, y) \land \Obj S_{\sigma}(x, y)) \lif {}]] \\
&\qquad (\Obj Q_{q_j}(x, y') \land \Obj S_{\sigma'}(x, y') \land
\varphi(x, y) \land \psi(y')))
\end{align*}
\end{enumerate}
येथे दिसते त्याप्रमाणे~$M$
च्या संक्रमणांचे वर्णन
करणारी वाक्ये~$\mathsf{T}(M,w)$ मधील
संबंधित वाक्यांसारखीच आहेत;
फक्त शेवटी $\psi(y')$
जोडले आहे. $\psi(y')$
मुळे “पुढच्या” स्थितिवर्णनाचा
क्रमांक $y'$ हा आधीच्या
सर्व क्रमांकांहून वेगळा
असतो: $\Obj 0$,
$\Obj 0'$, \dots.
























\end{enumerate}
$\mathsf{T}'(M, w)$ हे ट्यूरिंग
यंत्र~$M$ आणि आदान~$w$
यांच्यासाठी वरील सर्व
वाक्यांचे
संयुक्त विधान असू द्या.

\begin{lem}\label{tur:und:tra:lem:halts-sat}
  आदान~$w$ वर सुरू
  केलेले $M$ थांबत असेल,
  तर $\mathsf{T}'(M,w) \land \mathsf{E}(M,w)$
  ला सांत प्रतिरूप असते.
\end{lem}

\begin{proof}
  \ref{tur:und:ver:lem:halt-if-valid}
  च्या पुराव्यातील अनंत प्रतिरूपाच्या बांधणीप्रमाणे
  सांत रचना $\Struct{M'}$ बांधा;
  मात्र पुढील बदल करा:
  \begin{align*}
    \Domain{M'} & = \{0, \dots, n\},\\
    \Assign{\prime}{M'}(x) & =
    \begin{cases}
      x + 1 &\text{$x < n$ असेल तर}\\
      n &\text{अन्यथा,}
    \end{cases} \\
    \tuple{x,y} \in \Assign{<}{M'} &\text{$x < y$ किंवा $x = y = n$ असेल तेव्हा आणि केवळ तेव्हाच,}
  \end{align*}
  येथे $n = \max(k+1,\len{w}+1)$
  आणि~$k$ ही अशी
  लघुतम संख्या आहे
  की आदान~$w$ वर
  सुरू झालेले $M$
  हे~$k$ पायऱ्यांनंतर
  थांबले आहे.
  $\Sat{M'}{\mathsf{T}'(M,w) \land \mathsf{E}(M,w)}$
  याची पडताळणी सरावासाठी
  सोडतो.

\end{proof}

\begin{prob}
  $\Sat{M'}{\mathsf{T}'(M,w) \land \mathsf{E}(M,w)}$
  हे सिद्ध करून
  \ref{tur:und:tra:lem:halts-sat}
  चा पुरावा पूर्ण करा.
\end{prob}

\begin{lem}\label{tur:und:tra:lem:sat-halts}
  $\mathsf{T}'(M,w) \land \mathsf{E}(M,w)$
  ला सांत प्रतिरूप
  असेल, तर आदान~$w$
  वर सुरू झालेले
  $M$ थांबते.
\end{lem}

\begin{proof}
  प्रतिपरिवर्तन सिद्ध करू.
  समजा आदान~$w$
  वर सुरू झालेले
  $M$ थांबत नाही.
  $\mathsf{T}'(M,w) \land \mathsf{E}(M,w)$
  ला प्रतिरूपच नसेल,
  तर पुरावा पूर्ण.
  म्हणून~$\Struct{M}$
  हे~$\mathsf{T}'(M,w) \land \mathsf{E}(M, w)$
  चे प्रतिरूप आहे
  असे समजा. ते
  सांत असू शकत
  नाही हे दाखवू.


  \ref{tur:und:ver:lem:config}
  प्रमाणे, आदान~$w$
  वर सुरू केलेले~$M$
  हे~$n$ पायऱ्यांनंतर
  थांबलेले नसेल आणि
  किमान एक पायरी
  झालेली असेल, तर
  $\mathsf{T}'(M,w)
  \Entails \chi(M, w, n) \land \psi(\num{n})$
  हे सिद्ध करता येते.
  आदान~$w$ वर
  सुरू झालेले~$M$
  थांबत नसल्यामुळे,
  प्रत्येक सकारात्मक
  $n \in \Nat$ साठी
  $\mathsf{T}'(M,w) \Entails \chi(M, w, n) \land
  \psi(\num{n})$.
  \ref{tur:und:rep:prop:mlessk}
  नुसार सर्व~$k < n$
  साठी $\mathsf{T}'(M,w) \Entails \num{k} < \num{n}$
  हेही लक्षात घ्या.
  तसेच $\psi(\num{n})
  \Entails \num{k} < \num{n} \lif
  \eq/[\num{k}][\num{n}]$.
  म्हणून सर्व~$k < n$
  साठी $\Sat{M}{\eq/[\num{k}][\num{n}]}$;
  म्हणजे अनंत संख्यांक
  पदे~$\num{k}$ यांना
  $\Struct{M}$ मध्ये
  परस्परभिन्न मूल्ये
  असलीच पाहिजेत.
  त्यासाठी $\Domain{M}$
  अनंत असावे लागते;
  म्हणून~$\Struct{M}$
  हे~$\mathsf{T}'(M,w) \land \mathsf{E}(M, w)$
  चे सांत प्रतिरूप
  असू शकत नाही.

\end{proof}

\begin{prob}
  आदान~$w$ वर
  सुरू केलेले~$M$
  हे~$n$ पायऱ्यांनंतर
  थांबलेले नसेल आणि
  किमान एक पायरी
  झालेली असेल,
  तर $\mathsf{T}'(M,w) \Entails \psi(\num{n})$
  हे सिद्ध करून
  \ref{tur:und:tra:lem:sat-halts}
  चा पुरावा पूर्ण करा.
\end{prob}

\begin{thm}[ट्राख्टेनब्रोटचे प्रमेय]
  \label{tur:und:tra:thm:trakhtenbrodt}
  प्रथम-क्रम तर्कशास्त्रातील
  कोणत्याही वाक्याला सांत प्रतिरूप
  आहे का (म्हणजे ते
  सांत प्रतिरूपात पूर्ततायोग्य आहे का),
  हे ठरवणे अनिर्णेय आहे.
\end{thm}

\begin{proof}
  सांत प्रतिरूपात पूर्ततायोग्यता ठरवणारे
  ट्यूरिंग यंत्र~$F$
  आहे असे समजा.
  मग कोणतेही ट्यूरिंग
  यंत्र~$M$ आणि
  आदान~$w$ दिले असता
  $\mathsf{T}'(M,w) \land \mathsf{E}(M,w)$
  हे वाक्य संगणित करून
  त्याला सांत प्रतिरूप
  आहे का हे~$F$
  ने ठरवता येईल.
  \ref{tur:und:tra:lem:halts-sat} आणि \ref{tur:und:tra:lem:sat-halts}
  नुसार त्याला सांत
  प्रतिरूप असते तेव्हा आणि केवळ तेव्हाच
  आदान~$w$ वर
  सुरू झालेले $M$
  थांबते. म्हणून~$F$
  वापरून थांबण्याची
  समस्या सोडवता येईल;
  पण ती सोडवता येत
  नाही हे आपल्याला
  माहीत आहे.
\end{proof}

\begin{cor}\label{tur:und:tra:cor:fproof-incomp}
  \emph{सांत} वैधतेसाठी
  निर्दोष आणि संपूर्ण, तसेच निष्पत्तींना
  प्रभावीरीत्या प्रगणित करता येईल अशी
  कोणतीही निष्पत्ती-पद्धती असू शकत नाही.
  म्हणजे $\Proves \psi$ तेव्हा आणि केवळ तेव्हाच,
  जेव्हा प्रत्येक सांत रचना~$\Struct{M}$ मध्ये $\Sat{M}{\psi}$,
  अशी प्रभावी निष्पत्ती-पद्धती असू शकत नाही.

\end{cor}

\begin{proof}
  सराव.
\end{proof}

\begin{prob}
  \ref{tur:und:tra:cor:fproof-incomp}
  सिद्ध करा. प्रत्येक
  सांत रचनेत $\psi$ सत्य
  असेल तेव्हा आणि केवळ तेव्हाच
  $\lnot \psi$ सांत प्रतिरूपात
  पूर्ततायोग्य नसेल,
  हे लक्षात घ्या.
  \ref{tur:und:uns:thm:valid-ce}
  मधील अर्थाने सांत
  प्रतिरूपात पूर्ततायोग्यता अर्धनिर्णेय
  का आहे, हे स्पष्ट
  करा. सांत वैधतेसाठी
  निष्पत्ती प्रभावीरीत्या प्रगणित करता येणारी
  निष्पत्ती-पद्धती
  असेल, तर सांत
  प्रतिरूपात पूर्ततायोग्यता निर्णेय
  होईल, असा युक्तिवाद
  करण्यासाठी हे वापरा.
\end{prob}




\part{अपूर्णता}
